\chapter{The Stability of More General Flows Between Coaxial Cylinders}
\label{ch:8}

\section{Introduction}
\label{sec:75}

In the last chapter we considered the stability of the simple Couette flow between rotating cylinders. In this chapter we shall consider some more general flows between coaxial cylinders when, in addition to rotation, a constant transverse or axial pressure gradient is present. In the former case, the streamlines continue to be circular and satisfy the requirement of a Couette flow; but the flow differs from the one considered in Chapter~VII in that a Poiseuille flow in the transverse direction is superposed over the distribution of rotational velocities,
\begin{equation}
\label{eq:8-1}
\Omega = A + B/r^2,
\end{equation}
given in equation~VII\,(144). On the other hand, when a constant pressure gradient $\partial p/\partial z$ along the common axis of the cylinders is present, a Poiseuille flow in the axial direction is superposed over the same distribution of rotational velocities. The streamlines are then no longer circular and the flow cannot be classed under Couette flow. Indeed, as we shall see, the superposition of an axial flow over a rotational flow introduces certain essentially new elements into the problem.

\section{The stability of viscous flow in a curved channel}
\label{sec:76}

Consider equations~VII\,(134)--(138) in the case when the conditions are stationary, a constant pressure gradient $(\partial p/\partial\theta)_0$ is acting in the transverse direction, the cylinders are at rest, and only transverse motions are present. Under these circumstances, the equations of motion allow the stationary solution (cf.\ equations~VII\,(141) and (142))
\begin{equation}
\label{eq:8-2}
u_r = u_z = 0 \quad \text{and} \quad u_\theta = V(r),
\end{equation}
provided
\begin{equation}
\label{eq:8-3}
\frac{1}{\rho}\frac{\partial p}{\partial r} = \frac{V^2}{r},
\end{equation}
and
\begin{equation}
\label{eq:8-4}
\nu\!\left(\nabla^2 V - \frac{V}{r^2}\right) = \nu \frac{d}{dr}\!\left(\frac{1}{r}\frac{d}{dr}\,rV\right) = \frac{1}{\rho r}\!\left(\frac{\partial p}{\partial\theta}\right)_{\!0}.
\end{equation}
The general solution of equation~\eqref{eq:8-4} is
\begin{equation}
\label{eq:8-5}
V = \frac{1}{2\rho\nu}\!\left(\frac{\partial p}{\partial\theta}\right)_{\!0} r\log r + Cr + \frac{D}{r},
\end{equation}
where $C$ and $D$ are two constants of integration. If, as we have assumed, the flow takes place between two concentric cylinders (of radii $R_1$ and $R_2\,(> R_1)$) at rest, we must require that $V$ vanishes at $r = R_1$ and $R_2$. These latter conditions determine $C$ and $D$; and we find
\begin{equation}
\label{eq:8-6}
C = -\frac{1}{2\rho\nu}\!\left(\frac{\partial p}{\partial\theta}\right)_{\!0}\frac{R_2^2\log R_2 - R_1^2\log R_1}{R_2^2 - R_1^2}
\end{equation}
and
\begin{equation}
\label{eq:8-6b}
D = \frac{1}{2\rho\nu}\!\left(\frac{\partial p}{\partial\theta}\right)_{\!0}\frac{R_1^2 R_2^2}{R_2^2 - R_1^2}\log\frac{R_2}{R_1}.
\end{equation}

When the difference $d$ in the radii of the two cylinders is small compared with their mean radius $\tfrac{1}{2}(R_2+R_1)$, the velocity distribution given by equations~\eqref{eq:8-5} and \eqref{eq:8-6} reduces to the familiar form of the Poiseuille flow between parallel planes. Thus, letting
\begin{equation}
\label{eq:8-7}
\xi = (r - R_1)/(R_2 - R_1),
\end{equation}
we find that to order $(d/R_1)^2$ the velocity distribution approximates to
\begin{equation}
\label{eq:8-8}
V(\xi) = 6V_m\,\xi(1-\xi),
\end{equation}
where
\begin{equation}
\label{eq:8-9}
V_m = -\frac{d^3}{12\rho R_1}\!\left(\frac{\partial p}{\partial\theta}\right)_{\!0}
\end{equation}
is the mean velocity across the channel.

The stability of the flow represented by the foregoing solution of the equations of motion was first considered by Dean. More recently, it has been reconsidered by Reid using methods similar to those described in other connexions in the earlier chapters.

Before we enter into a full discussion of the stability of the flow~\eqref{eq:8-8}, we may note that Rayleigh's discriminant (defined in equation~VII\,(18)) is positive for $0 \leqslant \xi \leqslant \tfrac{1}{4}$ and negative for $\tfrac{1}{4} < \xi \leqslant 1$. Accordingly, in the absence of viscosity the flow is unstable; and the instability is principally to be associated with the violation of Rayleigh's criterion beyond $\xi = \tfrac{1}{4}$.

\subsection{The perturbation equations}
\label{sec:76a}

Equations~VII\,(161) and (162) are applicable to the problem on hand if we ascribe to $V$ in these equations its present meaning.

When the gap $d = R_2 - R_1$ between the cylinders is small compared to their mean radius $\tfrac{1}{2}(R_2 + R_1)$, the appropriate expression for $V$ is that given in equation~\eqref{eq:8-8}. In the same approximation, the perturbation equations are (cf.\ equations~VII\,(193) and (194))
\begin{equation}
\label{eq:8-10}
(D^2 - a^2 - \sigma)(D^2 - a^2)u = \frac{12V_m\,d^2}{R_1\,\nu}\,a^2\xi(1-\xi)v
\end{equation}
and
\begin{equation}
\label{eq:8-11}
(D^2 - a^2 - \sigma)v = \frac{6V_m\,d}{R_1\,\nu}(1-2\xi)u,
\end{equation}
where $\sigma$ and $a$ have the same meanings as in equation~VII\,(192). By the further transformation,
\begin{equation}
\label{eq:8-12}
\mathbf{u} \to \frac{12V_m\,d^2}{R_1^2\,\nu}\,a^2\mathbf{u},
\end{equation}
the equations take the forms:
\begin{equation}
\label{eq:8-13}
(D^2 - a^2 - \sigma)(D^2 - a^2)u = \xi(1-\xi)v
\end{equation}
and
\begin{equation}
\label{eq:8-14}
(D^2 - a^2 - \sigma)v = \Lambda a^2(1-2\xi)u,
\end{equation}
where
\begin{equation}
\label{eq:8-15}
\Lambda = \frac{72 V_m^2\,d^3}{R_1^3\,\nu^2} = 72\,\mathrm{R}^2\,\frac{d}{R_1}
\end{equation}
and
\begin{equation}
\label{eq:8-16}
\mathrm{R} = V_m d/\nu
\end{equation}
is the Reynolds number of the mean flow.

The boundary conditions with respect to which equations~\eqref{eq:8-13} and \eqref{eq:8-14} must be solved are (cf.\ equation~VII\,(200))
\begin{equation}
\label{eq:8-17}
\mathbf{u} = Du = v = 0 \quad \text{for } \xi = 0 \text{ and } 1.
\end{equation}

In view of the physical similarity of this problem with the problem considered in Chapter~VII, \S\,71, it is likely that in this case also the onset of instability will be as a stationary secondary flow. No theoretical justification has been given for this supposition; and in spite of experimental evidence favouring it (cf.\ \S\,77\,(c)), the matter deserves examination (cf.\ remarks on p.~317).

When the marginal state is stationary, the equations to be solved are
\begin{equation}
\label{eq:8-18}
(D^2 - a^2)^2 u = \xi(1-\xi)v
\end{equation}
and
\begin{equation}
\label{eq:8-19}
(D^2 - a^2)v = \Lambda a^2(1-2\xi)u,
\end{equation}
together with the boundary conditions~\eqref{eq:8-17}.

\subsection{The solution of the characteristic value problem for the case $\sigma = 0$}
\label{sec:76b}

The characteristic value problem presented by equations~\eqref{eq:8-17}--\eqref{eq:8-19} can be solved by the method described in \S\,71\,(a).

Since $v$ is required to vanish at $\xi = 0$ and 1, we expand it in a sine series of the form
\begin{equation}
\label{eq:8-20}
v = \sum_{n=1}^{\infty} C_n \sin m\pi\xi.
\end{equation}

With $v$ chosen in this manner, we next solve the equation,
\begin{equation}
\label{eq:8-21}
(D^2 - a^2)^2 u = \xi(1-\xi)\sum_{n=1}^{\infty} C_n \sin m\pi\xi,
\end{equation}
obtained by inserting \eqref{eq:8-20} in \eqref{eq:8-18}, and arrange that the solution satisfies the four remaining boundary conditions on $u$. With $u$ determined in this fashion and $v$ given by \eqref{eq:8-20}, equation~\eqref{eq:8-19} will lead to the required secular equation for~$\Lambda$.

The solution of equation~\eqref{eq:8-21} is straightforward and can be written in the form (cf.\ equation~VII\,(205))
\begin{multline}
\label{eq:8-22}
u = \sum_n \frac{C_n}{(m^2\pi^2 + a^2)^2}\biggl\{A_1^{(m)}\cosh a\xi + B_1^{(m)}\sinh a\xi^{\vphantom{X}} + A_2^{(m)}\xi\cosh a\xi + B_2^{(m)}\xi\sinh a\xi \\
+ \biggl[\xi(1-\xi) + \frac{4(5m^2\pi^2 - a^2)}{(m^2\pi^2 + a^2)^2}\biggr]\sin m\pi\xi \\
+ \frac{4m\pi}{(m^2\pi^2 + a^2)}(1-2\xi)\cos m\pi\xi\biggr\},
\end{multline}
where the constants of integration $A_1^{(m)}$, $A_2^{(m)}$, $B_1^{(m)}$, and $B_2^{(m)}$ are to be determined by the boundary conditions $\mathbf{u} = Du = 0$ at $\xi = 0$ and~1. These latter conditions lead to the equations:
\begin{equation}
\label{eq:8-23}
A_1^{(m)} = -\frac{4m\pi}{m^2\pi^2 + a^2},
\end{equation}
\begin{multline}
\label{eq:8-23b}
aB_1^{(m)} + A_2^{(m)} = -m\pi\biggl[\frac{4(5m^2\pi^2 - a^2)}{(m^2\pi^2 + a^2)^2} - \frac{8}{m^2\pi^2 + a^2}\biggr] = -\frac{12m\pi(m^2\pi^2 - a^2)}{(m^2\pi^2 + a^2)^2}, \\
\end{multline}
\begin{multline}
\label{eq:8-23c}
A_1^{(m)}\cosh a + B_1^{(m)}\sinh a + A_2^{(m)}\cosh a + B_2^{(m)}\sinh a = (-1)^{m+1}\frac{4m\pi}{(m^2\pi^2 + a^2)},
\end{multline}
\begin{multline}
\label{eq:8-23d}
A_1^{(m)}a\sinh a + B_1^{(m)}a\cosh a + A_2^{(m)}(\cosh a + a\sinh a) + {} \\
B_2^{(m)}(\sinh a + a\cosh a) = (-1)^{m+1}m\pi\frac{12(m^2\pi^2 - a^2)}{(m^2\pi^2 + a^2)^2}.
\end{multline}
On solving these equations we find (cf.\ equations~VII\,(207) and (208)):
\begin{align}
\label{eq:8-24}
A_1^{(\infty)} &= -\frac{4m\pi}{m^2\pi^2 + a^2}, \notag\\[4pt]
B_1^{(\infty)} &= +\frac{m\pi}{\Delta}\bigl(\alpha_m a + \beta_m(\sinh a + a\cosh a) - \gamma_m\sinh a\bigr), \notag\\[4pt]
A_2^{(\infty)} &= -\frac{m\pi}{\Delta}\bigl(\alpha_m\sinh^2 a + \beta_m\,a(\sinh a + a\cosh a) - \gamma_m\,a\sinh a\bigr), \notag\\[4pt]
B_2^{(\infty)} &= +\frac{m\pi}{\Delta}\bigl(\alpha_m(\sinh a\cosh a - a) + \beta_m\,a^2\sinh a - \gamma_m(a\cosh a - \sinh a)\bigr),
\end{align}
where
\begin{equation}
\label{eq:8-25}
\Delta = \sinh^2 a - a^2,
\end{equation}
\begin{equation}
\label{eq:8-25b}
\alpha_m = \frac{12(m^2\pi^2 - a^2)}{(m^2\pi^2 + a^2)^2},\quad
\beta_m = \frac{4}{m^2\pi^2 + a^2}\bigl[(-1)^m + \cosh a\bigr],
\end{equation}
and
\begin{equation}
\label{eq:8-25c}
\gamma_m = (-1)^{m+1}\alpha_m + \frac{4}{m^2\pi^2 + a^2}\,a\sinh a.
\end{equation}

Now substituting for $v$ and $u$ from equations~\eqref{eq:8-20} and \eqref{eq:8-22} in equation~\eqref{eq:8-19}, we obtain
\begin{multline}
\label{eq:8-26}
\sum_n C_n(n^2\pi^2 + a^2)\sin n\pi\xi \\
= \Lambda a^2(2\xi - 1)\sum_n \frac{C_n}{(m^2\pi^2 + a^2)^2}\biggl\{A_1^{(m)}\cosh a\xi + B_1^{(m)}\sinh a\xi + {} \\
A_2^{(m)}\xi\cosh a\xi + B_2^{(m)}\xi\sinh a\xi \\
+ \biggl[\xi(1-\xi) + \frac{4(5m^2\pi^2 - a^2)}{(m^2\pi^2 + a^2)^2}\biggr]\sin m\pi\xi \\
+ \frac{4m\pi}{(m^2\pi^2+a^2)}(1-2\xi)\cos m\pi\xi\biggr\}.
\end{multline}
Multiplying equation~\eqref{eq:8-26} by $\sin n\pi\xi$ and integrating over the range of~$\xi$, we obtain a system of linear homogeneous equations for the constants $C_n = C_n/(m^2\pi^2 + a^2)^2$; and the requirement that these constants do not all vanish leads to the secular equation
\begin{multline}
\label{eq:8-27}
\bigl|\bigl|(1-2\xi)\cosh a\xi|n\rangle A_1^{(m)} + (1-2\xi)\sinh a\xi|n\rangle B_1^{(m)} + {} \\
+ \langle\xi(- 2\xi^2)\cosh a\xi|n\rangle A_2^{(m)} + \langle(\xi-2\xi^2)\sinh a\xi|n\rangle B_2^{(m)} + X_{mn} + {} \\
+ \frac{4(5m^2\pi^2 - a^2)}{(m^2\pi^2+a^2)}\,Y_{mn} + \frac{4m\pi}{m^2\pi^2+a^2}\,\frac{8mn}{a^2}\,\Lambda_1\biggr|\biggr| = 0,
\end{multline}
where
\begin{equation}
\label{eq:8-28}
X_{mn} = \int_0^1 (1-2\xi)\xi(1-\xi)\sin m\pi\xi\sin n\pi\xi\,d\xi = \frac{12}{n^4}\frac{(m+n)^4 - (m-n)^4}{(m^2-n^2)^4} - \frac{4mn}{\pi^2(m^2-n^2)^2}
\end{equation}
if $n+m$ is odd and zero otherwise,
\begin{equation}
\label{eq:8-29}
Y_{mn} = \int_0^1 (1-2\xi)^2\cos m\pi\xi\sin n\pi\xi\,d\xi = \frac{2n}{n(n^2-m^2)} - \frac{8}{n^2}\frac{(n+m)^2 + (n-m)^2}{\pi^2(n^2-m^2)^2}
\end{equation}
if $n+m$ is odd and zero otherwise,
\begin{equation}
\label{eq:8-30}
\text{and}\quad Z_{mn} = \frac{8mn}{\pi^2(n^2-m^2)^2} \quad \text{if } n+m \text{ is odd and zero otherwise.}
\end{equation}
The various matrix elements which occur in \eqref{eq:8-27} are linear combinations of the elements:
\begin{align}
\label{eq:8-31}
\langle\cosh a\xi|n\rangle &= \frac{n\pi}{n^2\pi^2 + a^2}\bigl[1 + (-1)^{n+1}\cosh a\bigr], \notag\\[4pt]
\langle\sinh a\xi|n\rangle &= \frac{n\pi}{n^2\pi^2 + a^2}(-1)^{n+1}\sinh a, \notag\\[4pt]
\langle\xi\cosh a\xi|n\rangle &= (-1)^{n+1}\frac{n\pi}{n^2\pi^2 + a^2}\biggl[\cosh a - \frac{2a}{n^2\pi^2 + a^2}\sinh a\biggr], \notag\\[4pt]
\langle\xi\sinh a\xi|n\rangle &= \langle\sinh a\xi|n\rangle - \frac{2a}{n^2\pi^2 + a^2}\langle\cosh a\xi|n\rangle, \notag\\[4pt]
\langle\xi^2\cosh a\xi|n\rangle &= (-1)^{n+1}\frac{n\pi}{n^2\pi^2 + a^2}\biggl[\cosh a - \frac{4a}{n^2\pi^2 + a^2}\sinh a \notag\\
&\qquad - \frac{2(n^2\pi^2 - 3a^2)}{(n^2\pi^2 + a^2)^2}\langle\cosh a\xi|n\rangle\biggr], \notag\\[4pt]
\langle\xi^2\sinh a\xi|n\rangle &= (-1)^{n+1}\frac{n\pi}{n^2\pi^2 + a^2}\biggl[\sinh a - \frac{4a}{n^2\pi^2 + a^2}\cosh a\biggr] - {} \notag\\
&\qquad \frac{2(n^2\pi^2 - 3a^2)}{(n^2\pi^2+a^2)^2}\langle\sinh a\xi|n\rangle.
\end{align}

\subsection{Numerical results}
\label{sec:76c}

Reid has solved the secular equation~\eqref{eq:8-27} in the second and in the fourth approximations for a number of values of~$a$. The results of his calculations for values of $a$ in the range in which $\Lambda$ attains its minimum are summarized in Table~XXXVIII. From the results given in this table we infer that the Reynolds number R at which instability will set in and the wave number of the disturbance which will be manifested at onset are given by
\begin{equation}
\label{eq:8-32}
\mathrm{R}_c = 35.94\sqrt{\frac{R_1}{d}} \quad \text{and} \quad a_c = 3.96.
\end{equation}
The corresponding velocity profiles and cell pattern are illustrated in Figs.~88 and~89.

\begin{table}[ht]
\centering
\caption{Critical values of $\Lambda$ and related constants}
\label{tab:XXXVIII}
\begin{tabular}{cccc}
\hline
 & \multicolumn{2}{c}{$\Lambda$} & $\mathrm{R}(d/R_1)^{1/2}$ \\
$a$ & Second approx. & Fourth approx. & $= \sqrt{\Lambda/72}$ \\
\hline
3.90 & 91,591 & & \\
3.95 & 91,530 & 92,975 & 35.935 \\
4.00 & 91,565 & 93,001 & 35.940 \\
4.10 & 91,756 & & \\
\hline
\end{tabular}
\end{table}

A confirmation of the theoretical values~\eqref{eq:8-32} has recently been provided by the experiments of Brewster, Grosberg, and Nissan described in \S\,77\,(c) (see Fig.~96\,b, facing p.~359). They find
\begin{equation}
\label{eq:8-33}
\mathrm{R}_c(d/R_1) = 36.5 \pm 1.1 \quad \text{and} \quad a = 4.9 \pm 0.8 \text{ (measured),}
\end{equation}
in satisfactory agreement with the predicted values.

\figurePlaceholder{88}{The perturbation in the radial velocity (a) and the tangential velocity (b) at the onset of instability.}

It is of interest to compare the Reynolds number for the onset of instability given by \eqref{eq:8-33} with the Reynolds number at which instability sets in for strictly plane-parallel flows. It is known that in the latter case, the onset of instability occurs as overstable oscillations when\footnote{See, for example, C.~C.~Lin, \emph{The Theory of Hydrodynamic Stability}, p.~29, Cambridge, England, 1955. The factor $\frac{3}{4}$ allows for the fact that Lin's definition of the Reynolds number is in terms of the velocity $(= \frac{3}{2}V_m)$ at the centre of the channel and its half-width.}
\begin{equation}
\label{eq:8-34}
\mathrm{R} = \tfrac{3}{4}\times 5300 \simeq 7070.
\end{equation}
Clearly, the criterion given by \eqref{eq:8-32} must cease to be valid when
\begin{equation}
\label{eq:8-34b}
\mathrm{R} = 35.94\sqrt{\frac{R_1}{d}} > 7070,
\end{equation}
or
\begin{equation}
\label{eq:8-35}
R_1 > 3.86 \times 10^4 d.
\end{equation}

\figurePlaceholder{89}{The cell pattern at the onset of instability. The stream function, $\psi\;(\propto u\cos a\xi)$, has been normalized to unity and the cell pattern is drawn symmetrical about $z/d = 0$.}

Conversely, for the instability of plane-parallel flows to be observable as theoretically predicted, the channel in which the experiments are carried out must be so accurately parallel that the radius of curvature at all points exceeds $4\times 10^4 d$ by a sufficiently large margin.

The origin of our inability to obtain the result valid for strictly plane-parallel flows as the limit of our present result for $R_1/d \to \infty$ must lie, among other things, in the fact that we have not explored the possibility of overstable oscillations occurring. One may surmise that overstability is indeed possible; but that it normally occurs for a value of~$\Lambda$ in excess of that at which a stationary secondary flow manifests itself. However, the discovery of an overstable branch (if one such exists) will not, by itself, resolve the difficulty: the passage to the limit $R_1/d\to\infty$ clearly requires great care.

\section{The stability of viscous flow between rotating cylinders when a transverse pressure gradient is present}
\label{sec:77}

We pass on now to a consideration of the stability of a general Couette flow which consists of a superposition of a Poiseuille flow in the transverse direction (maintained by a pressure gradient) and a distribution of angular velocities (maintained by the rotation of the two cylinders). Such general Couette flows can be realized in an arrangement of the kind shown in Fig.~90: the cylinders are allowed to rotate independently of one another while a constant flow through the annulus is maintained by a suitable pumping circuit.

\figurePlaceholder{90}{Arrangement by which general Couette flows can be realized.}

The stationary flow whose stability we wish to consider is represented by a superposition of the solutions given in \S\S\,69 and 76; thus,
\begin{equation}
\label{eq:8-36}
\frac{V(r)}{r} = \frac{1}{r^2}\frac{1}{2\rho\nu}\!\left(\frac{\partial p}{\partial\theta}\right)_{\!0} r^2\log r + Cr^2 + D + Ar + \frac{B}{r},
\end{equation}
where $C$ and $D$ have the values given in equation~\eqref{eq:8-6} and $A$ and $B$ have the values given in equation~VII\,(145).

In the case when the gap width $d = R_2 - R_1$ is small compared to the mean radius $\tfrac{1}{2}(R_2+R_1)$, we have the approximate relation (cf.\ equations~VII\,(191) and equation~\eqref{eq:8-8})
\begin{equation}
\label{eq:8-37}
\frac{V(r)}{r} = \Omega_1\bigl[1-(1-\mu)\xi\bigr] + \frac{6V_m}{R_1}\xi(1-\xi).
\end{equation}
We shall restrict our discussion to this case of narrow gap widths.

\subsection{The perturbation equations for the case $(R_2-R_1) \ll \frac{1}{2}(R_2+R_1)$}
\label{sec:77a}

The relevant perturbation equations can be readily written down by combining the right-hand sides of equations~VII\,(193) and (194) and equations \S\,76 \eqref{eq:8-10} and \eqref{eq:8-11}. We thus obtain
\begin{equation}
\label{eq:8-38}
(D^2 - a^2 - \sigma)(D^2 - a^2)u = \frac{2\Omega_1 d^2}{\nu}a^2\bigl\{\Omega_1[1-(1-\mu)\xi] + \tfrac{6V_m}{R_1}\xi(1-\xi)\bigr\}v
\end{equation}
and
\begin{equation}
\label{eq:8-39}
(D^2 - a^2 - \sigma)v = \frac{2d^2}{\nu}\biggl\{A + \frac{3V_m}{R_1}\,d\,(1-2\xi)\biggr\}u.
\end{equation}
Let
\begin{equation}
\label{eq:8-40}
\lambda = \frac{6V_m}{R_1\Omega_1}.
\end{equation}
With the value of $A$ given in equation~VII\,(145),
\begin{equation}
\label{eq:8-41}
\frac{3V_m}{Ad} = -\frac{3V_m(1-\eta^2)}{d_1\,\eta^2(1-\mu/\eta^2)\,d} = -\frac{3V_m(1+\eta)}{R_1\,\Omega_1(1-\mu)}.
\end{equation}
For the case of narrow gaps under consideration, $\eta \sim 1$ and we can write
\begin{equation}
\label{eq:8-42}
\frac{3V_m}{Ad} = -\frac{6V_m}{R_1\Omega_1(1-\mu)} = -\frac{\lambda}{1-\mu}.
\end{equation}
With these definitions we can rewrite equations~\eqref{eq:8-38} and \eqref{eq:8-39} in the forms
\begin{equation}
\label{eq:8-43}
(D^2 - a^2 - \sigma)(D^2 - a^2)u = \frac{2\Omega_1\,d^2}{\nu}a^2\bigl\{1-(1-\mu)\xi + \lambda\xi(1-\xi)\bigr\}v
\end{equation}
and
\begin{equation}
\label{eq:8-44}
(D^2 - a^2 - \sigma)v = \frac{2Ad^2}{\nu}\biggl[1 - \frac{\lambda}{1-\mu}(1-2\xi)\biggr]u.
\end{equation}
By the further transformation
\begin{equation}
\label{eq:8-45}
u \to \frac{2\Omega_1\,d^2}{\nu}a^2 u,
\end{equation}
the equations become
\begin{equation}
\label{eq:8-46}
(D^2 - a^2 - \sigma)(D^2 - a^2)u = \bigl[1-(1-\mu)\xi + \lambda\xi(1-\xi)\bigr]v
\end{equation}
and
\begin{equation}
\label{eq:8-47}
(D^2 - a^2 - \sigma)v = -Ta^2\biggl[1 - \frac{\lambda}{1-\mu}(1-2\xi)\biggr]u,
\end{equation}
where
\begin{equation}
\label{eq:8-48}
T = -\frac{4A\Omega_1}{\nu^2}\,d^4
\end{equation}
is the Taylor number appropriate for narrow gaps (cf.\ equation~VII\,(198)).

Solutions of equations~\eqref{eq:8-46} and \eqref{eq:8-47} must be sought which satisfy the usual boundary conditions
\begin{equation}
\label{eq:8-49}
\mathbf{u} = Du = v = 0 \quad \text{for } \xi = 0 \text{ and } 1.
\end{equation}

\subsection{The solution of the characteristic value problem for the case $\sigma = 0$}
\label{sec:77b}

For the problem under consideration, the principle of the exchange of stabilities has not been established. Nevertheless, if on the basis of experimental evidence we assume that the onset of instability is as a stationary secondary flow, the equations to be solved are
\begin{equation}
\label{eq:8-50}
(D^2 - a^2)^2 u = \bigl[1-(1-\mu)\xi + \lambda\xi(1-\xi)\bigr]v
\end{equation}
and
\begin{equation}
\label{eq:8-51}
(D^2 - a^2)v = -Ta^2\biggl[1 - \frac{\lambda}{1-\mu}(1-2\xi)\biggr]u,
\end{equation}
together with the boundary conditions~\eqref{eq:8-49}.

The characteristic value problem presented by equations~\eqref{eq:8-49}--\eqref{eq:8-51} can be solved by the same method which has proved successful for the separate problems. The details of the analysis need not be repeated since the relevant equations can be readily written down by suitably combining the corresponding equations of \S\,71\,(a) and \S\,76\,(b).

The solution for the case $\mu = 0$ has been explicitly carried out by

\begin{table}[ht]
\centering
\caption{The critical Taylor numbers and the wave numbers of the associated disturbance for various values of $\lambda$}
\label{tab:XXXIX}
\begin{tabular}{cccccc}
\hline
$\lambda$ & $a_c$ & $T_c$ & $\Lambda$ & $a_c$ & $T_c$ \\
\hline
$30$ & 3.70 & $8.20\times 10^3$ & $-7.5$ & 5.00 & $2.33\times 10^3$ \\
$21$ & 3.70 & $1.40\times 10^4$ & $-7.75$ & 5.73 & $2.83\times 10^4$ \\
15 & 3.60 & $2.55\times 10^4$ & $-9.00$ & 6.35 & $3.46\times 10^4$ \\
$15$ & 3.45 & $4.55\times 10^4$ & $-15$ & & \\
$6$ & 3.30 & $8.30\times 10^4$ & $-7.50$ & 7.40 & $6.38\times 10^4$ \\
3 & 3.14 & $1.24\times 10^5$ & $-3.75$ & 5.30 & $6.00\times 10^5$ \\
1 & 3.13 & $1.55\times 10^5$ & $-4.50$ & 5.37 & \\
0.5 & 3.14 & $2.84\times 10^4$ & $-6.00$ & & $9.51\times 10^5$ \\
$0$ & 3.17 & $3.39\times 10^5$ & & & \\
$-0.5$ & 3.17 & $4.43\times 10^5$ & $-6.00$ & & \\
$-1.0$ & 3.14 & $5.44\times 10^5$ & $-8000$ & 4.70 & $3.43\times 10^6$ \\
$-1.5$ & 3.47 & $7.60\times 10^4$ & $-10{,}000$ & 4.35 & \\
$-2.0$ & 3.80 & $1.26\times 10^5$ & $-15{,}000$ & 4.35 & $6.92\times 10^6$ \\
\hline
\end{tabular}
\end{table}

Di~Prima. His results for the critical Taylor numbers and the wave numbers of the disturbance which will be manifested at marginal stability are given in Table~XXXIX. They are further illustrated in Figs.~91 and~92.

\subsection{The physical interpretation of the results}
\label{sec:77c}

The very peculiar dependence of $T_c$ and $a_c$ on $\lambda$, exhibited by the calculations, can be understood by examining how Rayleigh's criterion for the stability of inviscid rotational flow is violated in the fluid. The distribution of the transverse velocity in the fluid (for the case $\mu = 0$) is given by
\begin{equation}
\label{eq:8-52}
1 - \xi + \lambda\xi(1-\xi) = (1-\xi)(1+\lambda\xi);
\end{equation}
and Rayleigh's discriminant (equation~VII\,(18)) for the velocity distribution~\eqref{eq:8-52} is given by
\begin{equation}
\label{eq:8-53}
\Phi = (1-\xi)(1+\lambda\xi)(\lambda - 1 - 2\lambda\xi);
\end{equation}
and Rayleigh's criterion for stability is that $\Phi$ should be positive. From an examination of Fig.~93, it is clear that we must distinguish the following three cases:

\medskip
\noindent case~i: \quad $\lambda \geqslant 1$ when $\Phi \geqslant 0$ for $0 \leqslant \xi < \tfrac{1}{2}\bigl(1-\tfrac{1}{\lambda}\bigr)$,
\begin{equation}
\label{eq:8-54}
\text{and}\quad \Phi \leqslant 0\quad \text{for } \tfrac{1}{2}\bigl(1-\tfrac{1}{\lambda}\bigr) \leqslant \xi \leqslant 1;
\end{equation}

\noindent case~ii: \quad $+1 > \lambda \geqslant -1$ when $\Phi \leqslant 0$ for $0 \leqslant \xi \leqslant 1$;
\begin{equation}
\label{eq:8-55}
\end{equation}

\noindent and

\noindent case~iii: \quad $\lambda < -1$ when $\Phi \geqslant 0$ for $\frac{1}{|\lambda|} \leqslant \xi \leqslant \tfrac{1}{2}\bigl(1+\tfrac{1}{|\lambda|}\bigr)$,
\begin{equation}
\label{eq:8-56}
\text{and}\quad \Phi \leqslant 0\quad \text{for } 0 \leqslant \xi \leqslant \frac{1}{|\lambda|} \quad\text{and}\quad \tfrac{1}{2}\bigl(1+\tfrac{1}{|\lambda|}\bigr) \leqslant \xi \leqslant 1.
\end{equation}

The parts of the fluid which are stable or unstable according to Rayleigh's criterion are shown in Fig.~94. We observe that as $\lambda \to \pm\infty$, the parts of the fluid which are unstable become increasingly confined to the interval $\tfrac{1}{2} < \xi < 1$. (As we shall presently explain, the instability in the zone $0 \leqslant \xi \leqslant |\lambda|^{-1}$ plays no significant role for $\lambda \to -\infty$.) In the limit $\lambda\to\pm\infty$, the problem thus tends to the one which we have

\figurePlaceholder{91}{The variation of the critical Taylor number $T_c$ as a function of $\lambda\;(= 6V_m/R_1\Omega_1)$. The calculated theoretical relation is the curve labelled $A$ and the dashed line is a linear extrapolation of it. The data represent the points at which instability was observed to occur. Data to the left of the arrow are unduly influenced by small errors in $\mu$.}

\figurePlaceholder{92}{The variation of the critical wave number $a_c$ as a function of $\lambda$.}

considered in \S\,76. This is, indeed, evident from equations~\eqref{eq:8-50} and \eqref{eq:8-51} which, for $|\lambda|\to\infty$, tend to
\begin{equation}
\label{eq:8-57}
(D^2 - a^2)^2 u = \lambda\xi(1-\xi)v
\end{equation}
and
\begin{equation}
\label{eq:8-58}
(D^2 - a^2)v = \frac{T\lambda}{1-\mu}\,a^2(1-2\xi)u.
\end{equation}
By the further transformation $u \to \lambda u$, these equations become identical with equations~\eqref{eq:8-18} and \eqref{eq:8-19} of \S\,76 with the only difference that $\Lambda$ is now replaced by $T\lambda^2/(1-\mu)$. We conclude that for $|\lambda|\to\infty$, the critical Taylor number for the onset of instability tends asymptotically to the value given by (cf.\ Table~XXXVIII)
\begin{equation}
\label{eq:8-59}
\frac{T\lambda^2}{1-\mu} \to \Lambda_c = 9.300\times 10^4.\footnotemark
\end{equation}
\footnotetext{Note that the definition of $T$ includes, through $A$, a factor $(1-\mu)$.}

Therefore,
\begin{equation}
\label{eq:8-60}
T_c \to 9.300 \times 10^4\,\frac{1-\mu}{\lambda^2} \quad \text{as } |\lambda| \to \infty.
\end{equation}
The corresponding limiting value of $a_c$ is given by
\begin{equation}
\label{eq:8-61}
a_c \to 3.96 \quad \text{as } |\lambda| \to \infty.
\end{equation}
The asymptotic relations given by \eqref{eq:8-60} for $\lambda \to \pm\infty$ are included in

\figurePlaceholder{93}{The distribution of the transverse velocity $(1-\xi)(1+\lambda\xi)$ for various values of $\lambda$. The curves are labelled by the values of $\lambda$ to which they refer.}

\figurePlaceholder{94}{The parts of the fluid which are stable or unstable according to Rayleigh's criterion for the flow $(1-\xi)(1+\lambda\xi)$ as a function of $\lambda$.}

Fig.~91. The fact, that for $\lambda \to +\infty$ the asymptote is approached from below while it is approached from above for $\lambda\to-\infty$, can be understood when we observe that for large positive~$\lambda$, the interval in which the flow is unstable is in excess of one-half, while for large negative~$\lambda$, the wider of the two intervals in which the flow is unstable is less than one-half; and arguments similar to those used in \S\,71\,(e) lead one to infer that the critical Taylor number varies approximately as the inverse fourth power of the extent of the widest zone of instability, provided the different zones are well separated. These same arguments also account for the fact that the values of $a_c$ are less than the limiting value~3.96 for large positive~$\lambda$, while they are greater than~3.96 for large negative~$\lambda$.

Consider next the range $-1 \leqslant \lambda \leqslant +1$. In this range, Rayleigh's criterion is violated throughout the interval. Under these circumstances, we may, in a first approximation, replace the terms in equations~\eqref{eq:8-50} and \eqref{eq:8-51}, which allow for the variation of $V(\xi)$ through the interval, by their average values. The arguments here are the same as in \S\,71\,(d). With the replacements suggested, the equations become
\begin{equation}
\label{eq:8-62}
(D^2 - a^2)^2 u = \tfrac{1}{4}(1+\tfrac{1}{3}\lambda)v
\end{equation}
and
\begin{equation}
\label{eq:8-63}
(D^2 - a^2)v = -Ta^2 u.
\end{equation}
From these equations it is apparent that in the range $-1 \leqslant \lambda \leqslant +1$ we should expect the relations (cf.\ equation~VII\,(231))
\begin{equation}
\label{eq:8-64}
T_c = \frac{3416}{1+\lambda/3} \quad \text{and} \quad a_c = 3.1
\end{equation}
to hold approximately. This expectation is confirmed by the results of Di~Prima's exact calculations. Thus, for $\lambda = +1$ and $-1$, the approximate values of $T_c$ we deduce from \eqref{eq:8-64} are $2.56\times 10^3$ and $5.12\times 10^3$, respectively; and these values should be contrasted with the `exact' values $2.43\times 10^3$ and $5.42\times 10^5$.

We now turn to the origin of the sharp maximum of $T_c$ at $\lambda = -3.5$. This must clearly be traced to the relative extents of the two zones of instability which occur for $\lambda < -1$. The argument here is that if there are different zones of instability, well separated, the onset of instability will be determined, principally, by the zone of widest extent. For $-3 < \lambda < -1$, the wider of the two zones of instability adjoins the inner cylinder, while for $\lambda < -3$, it adjoins the outer cylinder; for $\lambda = -3$, the two zones of instability are of equal extent and they are separated by a stable zone also of the same extent. The minimum width of the largest extent zone of instability occurs, then, for $\lambda = -3$; and if our earlier argument (without the qualification `well separated') should be taken literally, the maximum Taylor number should occur for this value of $\lambda = -3$. But the unqualified argument should not be taken literally in this instance: for, the two zones of instability are \emph{not} well separated; and the amplitudes of the perturbations, which are expected to be large in the unstable zones, will not decrease sufficiently to evanescence in the intervening stable zone. We must attribute to this last circumstance the occurrence of the maximum Taylor number at the somewhat displaced value, $\lambda = -3.5$, where the larger extent of the outer stable zone is compensated by its wider separation from the inner unstable zone. The fact that the maximum value of $T_c$ should be associated with the maximum value of $a_c$ is also clear.

\subsection{Comparison with experimental results}
\label{sec:77d}

Experiments which verify the theoretical results of \S\,(b) have been carried out by Brewster, Grosberg, and Nissan. Their experimental arrangement is shown in Fig.~95. The outer cylinder $a$ (of inside diameter 13~in.) moulded in two semi-cylindrical sections from Perspex was transparent to allow visual observations. The inner cylinder $b$ (of

\figurePlaceholder{95}{A sketch of the experimental arrangement used by Brewster, Grosberg, and Nissan.}

outside diameter 12~in.) made of copper was perforated to enable the flow of liquid through it. The inner cylinder was rotated by means of the half-shafts~$c$; the outer cylinder was always at rest. A steady, controllable flow of liquid through the annulus~$d$ could be maintained by an external pump. The liquid flowed in and out of the apparatus through axial pipes $e$ and $f$ connected with stationary pressure and suction boxes $g$ and~$h$. A vertical meridional section of the annulus could be viewed through the tank~$i$ filled with the liquid. The onset of instability was viewed through this tank by transmitted light.

The experiments consisted in increasing either the rate of rotation or the rate of pumping until the onset of instability manifested itself by the formation of cells. The authors found that a small change in the speed of rotation (or the rate of pumping) usually resulted in the appearance or the disappearance of the cells.

The liquids used in the experiments were glycerine-water solutions. The flow through the annulus was visualized by one of three methods: by using the self-indicating properties of glycerine solutions; by introducing dyes; or, simply by observing the formation of bubbles at the

\figurePlaceholder{96}{The visualization of the onset of instability in general Couette flows by Brewster, Grosberg, and Nissan: (a) for the case of a pure pressure maintained flow; (b) for the case $\lambda\;(= 6V_m/R_1\Omega_1) = -3$, when the net flow is zero.}

onset of instability. The onset of instability visualized by two of these methods is shown in Fig.~96\,$a$ (by the first method) and Fig.~96\,$b$ (by the third method).

Fig.~96\,$b$ exhibits the onset of instability in the case when the cylinders are at rest: it occurs at the value of $\Lambda$ predicted in \S\,76 (see equation~\eqref{eq:8-33}). The parameter~$\lambda$ introduced in the theory is related to the net volume rate of flow~$\mathscr{Q}$ per unit length of the cylinder. The relation in question is (cf.\ equations~\eqref{eq:8-37} and \eqref{eq:8-40})
\begin{equation}
\label{eq:8-65}
\mathscr{Q} = (\tfrac{1}{2}R_1\Omega_1 + V_m)d = \tfrac{1}{2}R_1\Omega_1(1+\tfrac{1}{3}\lambda).
\end{equation}
Since $\mathscr{Q}$ and $\Omega_1$ are measured quantities, the value of $\lambda$ can be deduced. [Note $\mathscr{Q} = 0$ for $\lambda = -3$; for this reason Brewster, Grosberg, and Nissan call this the case of `completely reversed flow'; it is also the case for which the extents of the two unstable zones and of the intermediate stable zone are all equal.]

The experimental results of Brewster, Grosberg, and Nissan on the critical Taylor numbers are included in Fig.~91. It will be observed that the agreement with the theoretical calculations is satisfactory.

\section{The stability of inviscid flow between coaxial cylinders when an axial pressure gradient is present}
\label{sec:78}

The equations of inviscid flow (equations~VII\,(13)--(16)) allow the stationary solution
\begin{equation}
\label{eq:8-66}
u_r = 0, \quad u_\theta = V(r), \quad u_z = W(r),
\end{equation}
and
\begin{equation}
\label{eq:8-67}
p = \rho\int \frac{V^2(r)}{r}\,dr,
\end{equation}
where $V(r)$ and $W(r)$ are arbitrary functions of $r$.

In considering the stability of the flow represented by the foregoing solution of the equations of motion, we shall restrict ourselves to perturbations which are independent of~$\theta$. Let the perturbed state be described by
\begin{equation}
\label{eq:8-68}
u_r,\quad V + u_\theta,\quad W + u_z, \quad \text{and} \quad \varpi = \delta p/\rho,
\end{equation}
where, in accordance with our assumption, $u_r$, $u_\theta$, $u_z$, and $\varpi$ are functions only of $r$ and $z$. The linear equations governing the perturbation are
\begin{equation}
\label{eq:8-69}
\frac{\partial u_r}{\partial t} + W\frac{\partial u_r}{\partial z} - 2\frac{V}{r}\,u_\theta = -\frac{\partial\varpi}{\partial r},
\end{equation}
\begin{equation}
\label{eq:8-70}
\frac{\partial u_\theta}{\partial t} + W\frac{\partial u_\theta}{\partial z} + \left(\frac{dV}{dr} + \frac{V}{r}\right)u_r = 0,
\end{equation}
\begin{equation}
\label{eq:8-71}
\frac{\partial u_z}{\partial t} + W\frac{\partial u_z}{\partial z} + u_r\frac{dW}{dr} = -\frac{\partial\varpi}{\partial z},
\end{equation}
and the equation of continuity,
\begin{equation}
\label{eq:8-72}
\frac{\partial u_r}{\partial r} + \frac{u_r}{r} + \frac{\partial u_z}{\partial z} = 0.
\end{equation}

Following the standard procedure, we analyse the disturbance into normal modes and suppose that the various quantities describing the perturbation have a $(z,t)$-dependence given by
\begin{equation}
\label{eq:8-73}
e^{i(pt+kz)},
\end{equation}
where $p$ is a constant (which can be complex) and $k$ is the wave number of the disturbance in the $z$-direction.

Let $u(r)$, $v(r)$, $w(r)$, and $\varpi(r)$ denote the amplitudes of the perturbations $u_r$, $u_\theta$, $u_z$, and $\varpi$ whose $(z,t)$-dependence is given by \eqref{eq:8-73}. Equations~\eqref{eq:8-69}--\eqref{eq:8-72} then give
\begin{equation}
\label{eq:8-74}
i(p+kW)u - 2\Omega v = -D\varpi,
\end{equation}
\begin{equation}
\label{eq:8-75}
i(p+kW)v = -(D_*V)u,
\end{equation}
\begin{equation}
\label{eq:8-76}
i(p+kW)w + (DW)u = -ik\varpi,
\end{equation}
and
\begin{equation}
\label{eq:8-77}
D_*\,\mathbf{u} = -ikw,
\end{equation}
where
\begin{equation}
\label{eq:8-78}
D = d/dr \quad \text{and} \quad D_* = D + 1/r.
\end{equation}

Eliminating $w$ between equations~\eqref{eq:8-76} and \eqref{eq:8-77}, we find
\begin{equation}
\label{eq:8-79}
(p+kW)D_*\,\mathbf{u} - k(DW)u = ik^2\varpi.
\end{equation}
Differentiating this equation with respect to~$r$ and making use of equation~\eqref{eq:8-74}, we obtain
\begin{equation}
\label{eq:8-80}
D\bigl\{(p+kW)D_*\,\mathbf{u}\bigr\} - D\bigl\{k(DW)u\bigr\} = -ik^2\bigl[i(p+kW)u - 2\Omega v\bigr].
\end{equation}
On simplification and rearrangement, equation~\eqref{eq:8-80} reduces to the form
\begin{equation}
\label{eq:8-81}
(p+kW)(DD_*-k^2)u - k^2\Psi(r)u = 2ik^2\Omega v,
\end{equation}
where
\begin{equation}
\label{eq:8-82}
\Psi(r) = D^2W + (DW)(D-D_*) = D^2 W - \frac{1}{r}\frac{dW}{dr} = r\frac{d}{dr}\!\left(\frac{1}{r}\frac{dW}{dr}\right).
\end{equation}

Finally, substituting for $v$ from equation~\eqref{eq:8-75}, we obtain
\begin{equation}
\label{eq:8-83}
(p+kW)(DD_*-k^2)u - k^2\Psi(r)u = -k^2\Phi(r)\frac{u}{p+kW},
\end{equation}
where
\begin{equation}
\label{eq:8-84}
\Phi(r) = \frac{\Omega}{r}\frac{d}{dr}(r^2\Omega) = \frac{\Omega}{r}\frac{d}{dr}(rV) = \frac{\Omega^2}{r}\frac{d}{dr}(r^2\Omega)
\end{equation}
is Rayleigh's discriminant as previously defined.

The boundary conditions are
\begin{equation}
\label{eq:8-85}
\mathbf{u} = 0 \quad \text{for } r = R_1 \text{ and } R_2.
\end{equation}
where $R_1$ and $R_2\;(> R_1)$ are the radii of the two cylinders confining the fluid.

An alternative form of equation~\eqref{eq:8-83} which we shall find useful is obtained by the substitution
\begin{equation}
\label{eq:8-86}
i\chi = \frac{u}{p+kW}.
\end{equation}
This substitution is generally permissible only if $V \neq 0$; for, in this case (cf.\ equation~\eqref{eq:8-75})
\begin{equation}
\label{eq:8-87}
v = -(D_*V)\chi,
\end{equation}
and $\chi$ is required to be a regular function in the interval $(R_1, R_2)$. [Note that for the form of $V$ permissible under viscous flow, $D_*V = 2A$ = constant.] In terms of $\chi$, equations~\eqref{eq:8-74} and \eqref{eq:8-79} take the forms
\begin{equation}
\label{eq:8-88}
(p+kW)^2\chi = D\varpi
\end{equation}
and
\begin{equation}
\label{eq:8-89}
(p+kW)^2 D_*\chi = k^2\varpi.
\end{equation}
Now eliminating $\varpi$ between these equations, we obtain
\begin{equation}
\label{eq:8-90}
D\bigl\{(p+kW)^2 D_*\chi\bigr\} - k^2(p+kW)^2\chi = -k^2\Phi(r)\chi,
\end{equation}
and the appropriate boundary conditions are (cf.\ equation~\eqref{eq:8-86})
\begin{equation}
\label{eq:8-91}
\chi = 0 \quad \text{for } r = R_1 \text{ and } R_2.
\end{equation}

\subsection{The case of a pure axial flow}
\label{sec:78a}

When the initial stationary flow is a purely axial one, the transformation~\eqref{eq:8-86} is, in general, not permissible: it is permissible if $p$ is complex, or if $(p+kW)$ has no zero in the interval $(R_1,R_2)$; and in no other case.

We begin our considerations, then, with the equation,
\begin{equation}
\label{eq:8-92}
(p+kW)(DD_*-k^2)u - k^2\Psi u = 0.
\end{equation}
Equation~\eqref{eq:8-92}, which one obtains from equation~\eqref{eq:8-83} by setting $\Phi = 0$, has a singular point where $W = -p/k$. For stable oscillations (where $p$ is real), this singular point can occur in the interval $(R_1, R_2)$. However, if $p$ should be complex (as it will be for damped and overstable oscillations), the equation is regular for all real values of~$r$.

We shall now show that \emph{a necessary condition for the occurrence of overstable oscillations is that} $\Psi(r)$ \emph{changes sign in the interval} $(R_1, R_2)$. This theorem is the analogue of one due to Rayleigh that a necessary condition for the instability of an inviscid plane-parallel flow is that the velocity profile has a point of inflexion.

To prove the theorem, suppose that $p\;(= p_r + ip_i)$ is complex. Then we can rewrite equation~\eqref{eq:8-92} in the form
\begin{equation}
\label{eq:8-93}
\bigl(DD_* - k^2\bigr)u - \frac{k^2\Psi}{p+kW}\,u = 0.
\end{equation}
Multiplying this equation by $ru^*$ (where $u^*$ is the complex conjugate of~$u$) and integrating over the range of~$r$, we obtain (after an integration by parts)
\begin{equation}
\label{eq:8-94}
\int_{R_1}^{R_2} r\bigl(|D_*\,u|^2 + k^2|u|^2\bigr)\,dr + k\int_{R_1}^{R_2} \frac{r^2\Psi}{p+kW}\,r|u|^2\,dr = 0.
\end{equation}
(The integrated part vanishes on account of the boundary conditions.) The real and the imaginary parts of equation~\eqref{eq:8-94} must vanish separately. The vanishing of the imaginary part gives
\begin{equation}
\label{eq:8-95}
p_i\int_{R_1}^{R_2} \frac{\Psi}{|p+kW|^2}\,r|u|^2\,dr = 0.
\end{equation}

If $p_i \neq 0$ (as we have supposed), the integral must vanish; and for this to happen, a necessary condition is clearly that $\Psi(r)$ vanishes for some value of $r$ in $(R_1, R_2)$.

An alternative way of establishing the same result is instructive. Rewriting equation~\eqref{eq:8-93} in the form
\begin{equation}
\label{eq:8-96}
\frac{d}{dr}\!\left(r\frac{du}{dr}\right) - \frac{u}{r} - k^2 r u - \frac{k^2\Psi}{p+kW}\,ru = 0
\end{equation}
we subtract from the equation
\begin{equation}
\label{eq:8-97}
u^*\frac{d}{dr}\!\left(r\frac{du}{dr}\right) - \frac{|u|^2}{r} - k^2 r|u|^2 - \frac{k^2\Psi}{p+kW}\,r|u|^2 = 0
\end{equation}
its complex conjugate. We thus obtain
\begin{equation}
\label{eq:8-98}
\frac{d}{dr}\!\left[u^*r\frac{du}{dr} - u\!\left(r\frac{du^*}{dr}\right)\right] = -2ip_i\,\frac{k^2\Psi}{|p+kW|^2}\,|u|^2.
\end{equation}
Letting
\begin{equation}
\label{eq:8-99}
U = \tfrac{1}{2}ir\!\left(u^*\frac{du}{dr} - u\frac{du^*}{dr}\right)
\end{equation}
stand for a real quantity, we can rewrite equation~\eqref{eq:8-98} in the form
\begin{equation}
\label{eq:8-100}
\frac{dU}{dr} = p_i\frac{rk^2\Psi}{|p+kW|^2}\,|u|^2.
\end{equation}
Since $U = 0$ for $r = R_1$ and $R_2$ (in virtue of the boundary conditions on~$u$), $dU/dr$ must change its sign somewhere inside the interval $(R_1, R_2)$; and this requires that $\Psi(r)$ do the same.

A further result which can be established when $p$ is complex is that \emph{the real part of $-p$ must lie between the maximum and the minimum values of~$kW$}. This follows from equation~\eqref{eq:8-90},
\begin{equation}
\label{eq:8-101}
D\bigl\{(p+kW)^2 D_*\chi\bigr\} - k^2(p+kW)^2\chi = 0,
\end{equation}
which one obtains from equation~\eqref{eq:8-92} by the transformation~\eqref{eq:8-86}. (The transformation is legitimate since we have supposed that $p$ is complex.) By multiplying equation~\eqref{eq:8-101} by $r\chi^*$ and integrating over the range of~$r$, we obtain (after an integration by parts)
\begin{equation}
\label{eq:8-102}
\int_{R_1}^{R_2} (p+kW)^2r\bigl(|D_*\chi|^2 + k^2|\chi|^2\bigr)\,dr = 0.
\end{equation}
The imaginary part of this equation gives
\begin{equation}
\label{eq:8-103}
\int_{R_1}^{R_2} (p_r + kW)r\bigl(|D_*\chi|^2 + k^2|\chi|^2\bigr)\,dr = 0.
\end{equation}
Since $p_i \neq 0$ (by hypothesis), the integral must vanish and the stated limits on~$p_r$ follow.

We have seen that the vanishing of $\Psi$ somewhere inside the interval $(R_1, R_2)$ is a necessary condition for the instability of the basic flow; we shall now prove that for a wide class of velocity distributions this condition is also a sufficient one. The proof consists in showing that under the conditions a wave number $k_c$ exists for which the flow is stable and that it is unstable for the immediately neighbouring wave numbers.

Let
\begin{equation}
\label{eq:8-104}
\Psi(r) = rDD_*(W)r = 0 \quad \text{for } r = r_c \text{ where } W = W_c \text{ (say).}
\end{equation}
We shall suppose that
\begin{equation}
\label{eq:8-105}
K(r) = -\frac{\Psi(r)}{r(W-W_c)} = -\frac{DD_*(W)r}{W - W_c} > 0 \quad \text{throughout the interval.}
\end{equation}

Further restrictions which we shall impose on $W$ are:
\begin{equation}
\label{eq:8-106}
W \geqslant 0 \quad \text{in the interval,}
\end{equation}
and
\begin{equation}
\label{eq:8-106b}
W = 0 \quad \text{for } r = R_1 \text{ and } R_2.
\end{equation}
Consider, under these circumstances, the equation,
\begin{equation}
\label{eq:8-107}
DD_*\,\mathbf{u} + rK\mathbf{u} = k^2\mathbf{u},
\end{equation}
together with the boundary conditions $\mathbf{u} = 0$ for $r = R_1$ and $R_2$. This is a characteristic value problem of the classical Sturm--Liouville type; and from the general theory it follows that under the conditions stated, the characteristic values $k^2$ can be arranged in a monotonic decreasing sequence. The characteristic value problem, moreover, allows a variational formulation: its solution is equivalent to finding extremal values of the expression,
\begin{equation}
\label{eq:8-108}
k^2 = \frac{\displaystyle\int_{R_1}^{R_2} r\bigl\{rKu^2 - (D_*u)^2\bigr\}\,dr}{\displaystyle\int_{R_1}^{R_2} ru^2\,dr},
\end{equation}
which are stationary with respect to arbitrary variations of $\mathbf{u}$ subject only to the boundary conditions. Therefore, the largest of the characteristic values of $k^2$ represents the absolute maximum which the quantity on the right-hand side of \eqref{eq:8-108} can attain.

We shall now show that for a suitably chosen~$u$, $k^2$ given by equation~\eqref{eq:8-108} can assume a positive value; in view of what has been said, the establishment of this fact will guarantee the existence of a positive characteristic value for~$k^2$. The particular choice of $\mathbf{u}$ we make is $W/r$. With this choice, the numerator in the expression for $k^2$ becomes
\begin{equation}
\label{eq:8-109}
\int_{R_1}^{R_2}\bigl\{KW^2 - rD_*(W/r)D_*(W/r)\bigr\}\,dr.
\end{equation}
After an integration by parts, we are left with
\begin{equation}
\label{eq:8-110}
\int_{R_1}^{R_2} W\bigl[KW + DD_*(W/r)\bigr]\,dr,
\end{equation}
since the integrated part vanishes on account of the conditions~\eqref{eq:8-106} imposed on~$W$. Now according to equation~\eqref{eq:8-105},
\begin{equation}
\label{eq:8-111}
DD_*(W/r) = -K(W - W_c).
\end{equation}
Inserting this expression in \eqref{eq:8-110}, we obtain (cf.\ equations~\eqref{eq:8-105} and \eqref{eq:8-106})
\begin{equation}
\label{eq:8-112}
W_c\int_{R_1}^{R_2} WK\,dr > 0.
\end{equation}
Thus, with the choice of $\mathbf{u}$ made, $k^2$ is indeed positive, and we infer that a positive characteristic value for $k^2$ exists. Let $k_c^2$ denote this value; and let $\mathbf{u}_c$ be the proper function belonging to it. Then
\begin{equation}
\label{eq:8-113}
DD_*\,\mathbf{u}_c - \frac{\Psi(r)}{W - W_c}\,\mathbf{u}_c = k_c^2\,\mathbf{u}_c.
\end{equation}

Now rewrite equation~\eqref{eq:8-92} in the form
\begin{equation}
\label{eq:8-114}
DD_*\,\mathbf{u} - \frac{\Psi(r)}{W-c}\,u = k^2 u,
\end{equation}
where
\begin{equation}
\label{eq:8-115}
c = -p/k = c_r + ic_i \quad (\text{say})
\end{equation}
is allowed to be complex. We shall suppose that $k^2$ is an analytic function of the complex variable $c$ and consider its behaviour near~$k_c^2$. From equations~\eqref{eq:8-113} and \eqref{eq:8-114},
\begin{equation}
\label{eq:8-116}
\frac{d}{dr}\!\left(r\mathbf{u}_c\frac{du}{dr} - ru\frac{du_c}{dr}\right) - r\Psi u u_c\!\left(\frac{1}{W-c} - \frac{1}{W-W_c}\right) = (k^2 - k_c^2)ru\,u_c.
\end{equation}
Integrating this equation over the range of $r$ and remembering that both $\mathbf{u}$ and $\mathbf{u}_c$ vanish at the limits, we get
\begin{equation}
\label{eq:8-117}
(k^2 - k_c^2)\int_{R_1}^{R_2} ru_c\,u\,dr = -(c - W_c)\int_{R_1}^{R_2}\frac{r^2\Psi u_c}{(W-c)(W-W_c)}\,r|u|^2\,dr.
\end{equation}
We now let $k^2 \to k_c^2$, $c \to W_c$ and $u \to u_c$. In this limit, equation~\eqref{eq:8-117} becomes
\begin{equation}
\label{eq:8-118}
\left(\frac{dc}{dk^2}\right)_{k^2\to k_c^2}\int_{R_1}^{R_2} ru_c^2\,dr = \lim_{c\to W_c,\sigma\to 0}\int_{R_1}^{R_2}\frac{Kr^2 u_c^2}{(W-c_r)^2 + c_i^2}\,dr,
\end{equation}
where we have reintroduced the function $K$ defined in equation~\eqref{eq:8-105}. Writing explicitly the real and the imaginary parts of the integral on the right-hand side of equation~\eqref{eq:8-118}, we have
\begin{equation}
\label{eq:8-119}
\int_{R_1}^{R_2}\frac{Kr^2 u_c^2}{(W-c_r)^2 + c_i^2}\,dr = \int\frac{K(W-c_r)}{(W-c_r)^2+c_i^2}\,r^2 u_c^2\,dr + i\int\frac{c_i K}{(W-c_r)^2+c_i^2}\,r^2 u_c^2\,dr.
\end{equation}
In the limit $c_i\to 0$ and $c_r \to W_c$ the real part in \eqref{eq:8-119} tends to the principal value of the integral
\begin{equation}
\label{eq:8-120}
\int_{R_1}^{R_2}\frac{Kr^2 u_c^2}{W - W_c}\,dr,
\end{equation}
while the imaginary part tends to
\begin{multline}
\label{eq:8-121}
\lim_{c_i\to 0}\int_{R_1}^{R_2}\frac{c_i K(r)r^2 u_c^2}{(W-W_c)^2+c_i^2}\,dr = \bigl[K(r)r^2 u_c^2\bigr]_{r=r_c}\lim_{c_i\to 0}\int\frac{c_i\,dW}{(W-W_c)^2+c_i^2} \\
\shoveleft{\quad = \pi\biggl[\frac{K(r)r^2 u_c^2}{|dW/dr|}\biggr]_{r=r_c}}\\
= \pm\pi\biggl[\frac{K(r)r^2 u_c^2}{|dW/dr|}\biggr]_{r=r_c}\quad (c_i\to\pm 0).
\end{multline}
Combining the foregoing results, we have an equation of the form
\begin{equation}
\label{eq:8-122}
\left(\frac{dc}{dk^2}\right)_{k^2\to k_c^2} = A \pm iB \quad \text{where } B > 0.
\end{equation}
Alternatively, we can write
\begin{equation}
\label{eq:8-123}
\left(\frac{dc}{dk^2}\right)_{k^2\to k_c^2} = \frac{A \mp iB}{A^2+B^2}\quad (B > 0).
\end{equation}

From equation~\eqref{eq:8-123} it follows that for values of $k^2$ (on the real axis) different from~$k_c^2$ the imaginary part of $c$ is finite. For these same neighbouring wave numbers, the imaginary part of $p$ is finite; and this implies the instability of these modes.

We shall now return to consider the case when $\Psi(r)$ is of the same sign throughout. Then, from equation~\eqref{eq:8-95} it follows that $p_i = 0$; this means that $p$ cannot be complex; and this in turn implies that the flow is stable. However, in all these cases, the proper functions will be characterized by discontinuities in their derivatives. This is a consequence of the fact (which we shall verify directly) that \emph{any value of} $-W/k$ \emph{is an admissible solution for~$p$}. Suppose, for example, that for a particular~$p$, $p+kW$ vanishes only once in $(R_1, R_2)$. We may then start integrating equation~\eqref{eq:8-92} from either end and with arbitrary initial values for $du/dr$. The integrations can be continued inward until we come to the point where $p+kW = 0$. At this point the two values of $\mathbf{u}$ and $du/dr$ will presumably disagree. By a suitable choice of the relative initial values of $du/dr$, $\mathbf{u}$ can be made continuous; a discontinuity in $du/dr$ will, however, remain. But the occurrence of such a discontinuity is not inconsistent with equation~\eqref{eq:8-92}.

It appears, then, that any value of $-kW$ is a possible value of $p$. Are other real values admissible? If so, $(p+kW)$ will be of one sign throughout. And this we can exclude. For, if $(p+kW)$ does not have a zero in the interval $(R_1, R_2)$, the transformation~\eqref{eq:8-86} is a legitimate one and equation~\eqref{eq:8-101} is a valid equation to consider. From this latter equation we can deduce (cf.\ equation~\eqref{eq:8-102})
\begin{equation}
\label{eq:8-124}
\int_{R_1}^{R_2} (p+kW)^2 r\bigl(|D_*\chi|^2 + r^2\chi^2\bigr)\,dr = 0,
\end{equation}
an equation which cannot be satisfied by any real~$p$. Thus, the assumption that the transformation~\eqref{eq:8-86} is a legitimate one has led to a contradiction. We conclude that \emph{any value of $p$ not included between $-kW_{\max}$ and $-kW_{\min}$ is not admissible as a solution}.

\subsection{The general case when rotation is also present}
\label{sec:78b}

Returning to the general case, we first observe that the presence of rotation affects some aspects of the problem in a fundamental way. Thus, the existence of the relation (equation~\eqref{eq:8-75})
\begin{equation}
\label{eq:8-125}
i(p+kW)v = -(D_*V)u
\end{equation}
requires that $(p+kW)$ has no zero in the allowed range of $r$. Therefore, if $p$ is real, it cannot lie between $-kW_{\max}$ and $-kW_{\min}$; and this result is exactly the opposite of that which obtains in the case of a pure axial flow. This circumstance already suggests that rotation alters the character of the problem in a qualitative way and that the results for the case of a pure axial flow cannot be obtained by a simple passage to the limit $\Omega \to 0$. Indeed, the principal conclusions to which we shall arrive are: \emph{the instabilities associated with a pure axial flow cease to be relevant when rotation is present; and stability is determined exclusively by the distribution of the transverse velocities, i.e.\ by Rayleigh's criterion}.

The equation to be considered is (cf.\ equation~\eqref{eq:8-90})
\begin{equation}
\label{eq:8-126}
D\bigl\{(c-W)^2 D_*\chi\bigr\} - k^2(c-W)^2\chi = -\Phi(r)\chi,
\end{equation}
where $c = -p/k$; and the boundary conditions are
\begin{equation}
\label{eq:8-127}
\chi = 0 \quad \text{for } r = R_1 \text{ and } R_2.
\end{equation}

And the principal question concerns the reality, or otherwise, of~$c$.

The natural and the most direct way of looking at the problem presented by equations~\eqref{eq:8-126} and the boundary conditions~\eqref{eq:8-127} is to consider it as a characteristic value problem for $c$ for an assigned $k$ and given $W(r)$ and $\Phi(r)$. We shall presently see that considered in this way, the problem allows a variational formulation even though the characteristic value parameter $c$ enters the problem non-linearly.

There is a less direct, alternative way of looking at the problem which is also useful. By writing
\begin{equation}
\label{eq:8-128}
\hat{\Omega}(r) = \lambda\omega(r),
\end{equation}
where $\lambda$ is some chosen unit of angular velocity (which might be taken to be the angular velocity at $r = R_1$ (say)), we can consider the problem as a characteristic value problem for~$\lambda^2$. Thus, letting
\begin{equation}
\label{eq:8-129}
\phi(r) = \frac{2\omega}{r}\frac{d}{dr}(r^2\omega),
\end{equation}
we have
\begin{equation}
\label{eq:8-130}
D\bigl\{(c-W)^2 D_*\chi\bigr\} - k^2(c-W)^2\chi = -\lambda^2\phi(r)\chi.
\end{equation}
The advantage of this latter formulation is that the characteristic value parameter $\lambda^2$ enters the problem linearly; and, moreover, with some supplementary conditions, the known theorems of the classical Sturmian theory can be applied. However, it is important to remember that the solution of the physical problem requires that \emph{every} positive $\lambda^2 > 0$ be derived as a characteristic value for a given $\phi(r)$ and an assigned $k$ for suitably selected values of $c$ (real or complex).

\subsubsection{A variational principle for $c$}
\label{sec:78b-i}

Consider equations~\eqref{eq:8-126} and \eqref{eq:8-127} as presenting a characteristic value problem for~$c$. By multiplying equation~\eqref{eq:8-126} by $r\chi$ and integrating over the range of~$r$, we obtain, after an integration by parts,
\begin{equation}
\label{eq:8-131}
\int_{R_1}^{R_2}(c-W)^2r\bigl\{(D_*\chi)^2 + k^2\chi^2\bigr\}\,r\,dr = \int_{R_1}^{R_2}\Phi(r)r\chi^2\,dr.
\end{equation}
Letting
\begin{equation}
\label{eq:8-132}
I_1 = \int_{R_1}^{R_2}\bigl\{(D_*\chi)^2 + k^2\chi^2\bigr\}\,r\,dr,\quad I_2 = \int_{R_1}^{R_2} W\bigl\{(D_*\chi)^2 + k^2\chi^2\bigr\}\,r\,dr,
\end{equation}
\begin{equation}
\label{eq:8-132b}
I_3 = \int_{R_1}^{R_2} W^2\bigl\{(D_*\chi)^2 + k^2\chi^2\bigr\}\,r\,dr, \quad\text{and}\quad I_4 = \int_{R_1}^{R_2} \Phi(r)\chi r\,dr,
\end{equation}
we can rewrite equation~\eqref{eq:8-131} in the form
\begin{equation}
\label{eq:8-133}
c^2 I_1 - 2cI_2 + I_3 - I_4 = 0;
\end{equation}
or solving for $c$, we have
\begin{equation}
\label{eq:8-134}
c = \frac{1}{I_1}\bigl(I_2 \pm \sqrt{I_2^2 - I_1 I_3 + I_1 I_4}\bigr).
\end{equation}

Equation~\eqref{eq:8-133} provides the basis for a variational principle; to see this, consider the effect on $c$ (determined in accordance with equation~\eqref{eq:8-134}) of an arbitrary variation $\delta\chi$ in $\chi$ compatible only with the boundary conditions on~$\chi$. To the first order in the variation, we have
\begin{equation}
\label{eq:8-135}
2(cI_1 - I_2)\delta c = -(c^2\delta I_1 - 2c\delta I_2 + \delta I_3 - \delta I_4),
\end{equation}
where $\delta I_1$, $\delta I_2$, $\delta I_3$, and $\delta I_4$ are the corresponding variations in the integrals $I_1$, $I_2$, $I_3$, and $I_4$. After an integration by parts (in each case), we find that the variations $\delta I_1$, $\delta I_2$, $\delta I_3$, and $\delta I_4$ are given by:
\begin{align}
\label{eq:8-136}
\delta I_1 &= -2\int_{R_1}^{R_2}\delta\chi\bigl[DD_*\chi - k^2\chi\bigr]\,r\,dr, \notag\\[4pt]
\delta I_2 &= -2\int_{R_1}^{R_2}\delta\chi\bigl[D(WD_*\chi) - k^2 W\chi\bigr]\,r\,dr, \notag\\[4pt]
\delta I_3 &= -2\int_{R_1}^{R_2}\delta\chi\bigl[D(W^2 D_*\chi) - k^2 W^2\chi\bigr]\,r\,dr;
\end{align}
also,
\begin{equation}
\label{eq:8-137}
\delta I_4 = 2\int_{R_1}^{R_2}\delta\chi\,\Phi(r)\chi\,r\,dr.
\end{equation}
Inserting the foregoing expressions in \eqref{eq:8-135}, we obtain
\begin{equation}
\label{eq:8-138}
(cI_1-I_2)\delta c = \int_{R_1}^{R_2}\delta\chi\bigl\{D\bigl[(c-W)^2 D_*\chi\bigr] - k^2(c-W)^2\chi + \Phi(r)\chi\bigr\}\,r\,dr.
\end{equation}
We observe that the quantity which appears as a factor of $r\,\delta\chi$ under the integral sign on the right-hand side of equation~\eqref{eq:8-138} vanishes if equation~\eqref{eq:8-126} governing $\chi$ is satisfied.$\dagger$ Hence, a necessary and sufficient condition for $\delta c$ to vanish identically to the first order, for all small variations in $\chi$ subject only to the boundary conditions, is that $\chi$ be a solution of the characteristic value problem. While this theorem provides the basis for a variational treatment of the problem, the principal conclusion we wish to draw from it and equation~\eqref{eq:8-134} is that \emph{other things being equal, there are two, and only two, branches to the dispersion relation}.

\subsubsection{A variational principle for $\lambda^2$}
\label{sec:78b-ii}

Now consider equations~\eqref{eq:8-127} and \eqref{eq:8-130} as presenting a characteristic value problem for~$\lambda^2$. If we restrict ourselves to real values of $c$ not included between $W_{\max}$ and $W_{\min}$, we can apply the standard theorems of the classical Sturmian theory (cf.\ \S\,67\,(b)) and conclude: the characteristic values of $\lambda^2$ are all positive if $\phi(r)$ is everywhere positive and they are all negative if $\phi(r)$ is everywhere negative; but if $\phi(r)$ should change sign anywhere inside the interval $(R_1, R_2)$, then there are two sets of real characteristic values which have the limit points $\pm\infty$ and $-\infty$. Moreover, solving equation~\eqref{eq:8-130} with the boundary conditions~\eqref{eq:8-127} is equivalent to finding extremal values of $\lambda^2$ given by
\begin{equation}
\label{eq:8-139}
\lambda^2 = \frac{\displaystyle\int_{R_1}^{R_2}(c-W)^2r\bigl\{(D_*\chi)^2 + k^2\chi^2\bigr\}\,r\,dr}{\displaystyle\int_{R_1}^{R_2}\phi(r)r\chi^2\,dr},
\end{equation}
which are stationary with respect to arbitrary small variations of $\chi$ subject only to the boundary conditions~\eqref{eq:8-127}.

\subsubsection{The criterion for stability}
\label{sec:78b-iii}

The theorems enunciated in the preceding subsections enable us to derive a necessary and sufficient condition for stability.

Consider first the case when $\phi(r)$ is everywhere negative. Then, for every real $k$ and for all real values of $c$ (not included between $W_{\max}$ and $W_{\min}$), the characteristic values of $\lambda^2$ are all negative. Consequently, if $\lambda^2$ is to be positive, $c$ must be complex; and this means instability.

Consider next the case when $\phi(r)$ is everywhere positive. Then, for a given $k$ and $c$ (both assumed real and $c > W_{\max}$ or $c < W_{\min}$), the characteristic values, $\lambda_n^2$, of $\lambda^2$ are all positive. They can be arranged as an infinite, monotonic, increasing sequence; if $\lambda_n^2\;(n=1,2,\ldots)$ represents this arrangement, the proper function $\chi_n$ belonging to $\lambda_n^2$ will have exactly $(n-1)$ nodes in the interval $(R_1, R_2)$. To emphasize the dependence of $\lambda_n^2$ and $\chi_n$ on $c$ (for assigned~$k$) we shall write $\lambda_n^2(c)$ and $\chi_n(c)$.

Now by the variational principle of \S\,(ii), the smallest of the characteristic values $\lambda_1^2(c)$ represents the absolute minimum which the quantity on the right-hand side of equation~\eqref{eq:8-139} can attain. Therefore, if $\lambda^2(c;\,\chi)$ denotes the value of $\lambda^2$ given by \eqref{eq:8-139} for some assigned $c$ and chosen~$\chi$, then
\begin{equation}
\label{eq:8-140}
\lambda^2(c;\,\chi) \geqslant \lambda_1^2(c).
\end{equation}
The equality sign in \eqref{eq:8-140} can hold if and only if $\chi = \chi_1(c)$, i.e.\
\begin{equation}
\label{eq:8-141}
\lambda^2(c;\,\chi_1(c)) = \lambda_1^2(c).
\end{equation}

We shall now show that $\lambda_1^2(c)$ is a \emph{monotonic increasing function for} $c > W_{\max}$ \emph{and a monotonic decreasing function for} $c < W_{\min}$.

We first observe that according to equation~\eqref{eq:8-139}
\begin{equation}
\label{eq:8-142}
\lambda^2(c_2;\,\chi_1(c_1)) < \lambda^2(c_1;\,\chi_1(c_1)) = \lambda_1^2(c_1)\quad \text{if } W_{\max} < c_2 < c_1.
\end{equation}
On the other hand, by \eqref{eq:8-140},
\begin{equation}
\label{eq:8-143}
\lambda^2(c_2;\,\chi_1(c_1)) > \lambda_1^2(c_2).
\end{equation}
Hence, combining \eqref{eq:8-142} and \eqref{eq:8-143}, we have
\begin{equation}
\label{eq:8-144}
\lambda_1^2(c_2) < \lambda_1^2(c_1)\quad \text{if } W_{\max} < c_2 < c_1.
\end{equation}
In exactly the same way, we can show that
\begin{equation}
\label{eq:8-145}
\lambda_1^2(c_2) > \lambda_1^2(c_1)\quad \text{if } c_2 < c_1 < W_{\min}.
\end{equation}

It is also clear from \eqref{eq:8-139} that $\lambda^2(c)$ cannot be bounded as $c\to \pm\infty$.
Hence,
\begin{equation}
\label{eq:8-146}
\lambda_1^2(c) \to \infty \quad \text{as } c \to \pm\infty.
\end{equation}
Further, by letting $c$ approach $W_{\max}$ from above, or $W_{\min}$ from below, we can make $\lambda^2$ given by \eqref{eq:8-139} as small as we please; consequently,
\begin{equation}
\label{eq:8-147}
\lambda_1^2(c) \to 0 \quad \text{as } c \to W_{\max}+0 \quad \text{or } W_{\min}-0.
\end{equation}
(Note that $c = W_{\max}$ and $c = W_{\min}$ are excluded.)

From the foregoing results, it follows that for every assigned $k^2$ and $\lambda^2 > 0$, there exist \emph{two} real characteristic values for~$c$, one larger than $W_{\max}$ and one less than $W_{\min}$; and that the proper functions belonging to these two values have no nodes in $(R_1, R_2)$. By similarly considering characteristic values $\lambda_n^2$ (whose proper functions have $n-1$ nodes), we can deduce that for every assigned $k^2$ and $\lambda^2 > 0$ there are two real characteristic values for $c$ (one larger than $W_{\max}$ and one less than $W_{\min}$) whose proper functions have $n-1$ nodes. On the other hand, by the variational principle of \S\,(i), there can be, in this case, no more than two branches of the dispersion relation. And since we have already accounted for two branches with real $c$'s, there is no room for further complex values. The flow is, therefore, stable when $\phi(r)$ is everywhere positive.

When $\phi(r)$ changes sign in $(R_1, R_2)$, then the possible characteristic values of $\lambda^2$ (for any given $\lambda^2$ and all real $c$'s not included between $W_{\max}$ and $W_{\min}$) form two distinct sequences with limit points at $+\infty$ and $-\infty$. The existence of this latter sequence (with the limit point at $-\infty$) guarantees that for any given $\lambda^2 > 0$ there exist axial modes with complex characteristic values $c$ which ensure instability.

Thus, Rayleigh's criterion for the stability of a pure rotational flow continues to be valid in the presence of an arbitrary axial flow. This universality of Rayleigh's criterion is somewhat unexpected; but its origin is rooted in equation~\eqref{eq:8-75} which couples the radial and the transverse perturbations in the velocity. The present problem is, however, not unique: we have already encountered an example in Chapter~III (\S\,27\,(a), p.~97) in which the presence of rotation alters the character of the phenomenon equally.

\section{The stability of viscous flow between rotating coaxial cylinders when an axial pressure gradient is present}
\label{sec:79}

Consider equations~VII\,(134)--(136) in case the conditions are stationary, a constant pressure gradient $(\partial p/\partial z)_0$ is acting in the $z$-direction, the cylinders are in rotation, and no radial motions are present. Under these circumstances, the equations of motion allow the stationary solution
\begin{equation}
\label{eq:8-148}
u_r = 0, \quad u_\theta = V(r), \quad \text{and} \quad u_z = W(r),
\end{equation}
provided
\begin{equation}
\label{eq:8-149}
\frac{1}{\rho}\frac{dp}{dr} = \frac{V^2}{r},
\end{equation}
\begin{equation}
\label{eq:8-150}
\nu\!\left(\nabla^2 V - \frac{V}{r^2}\right) = \nu DD_*\,V = 0,
\end{equation}
and
\begin{equation}
\label{eq:8-151}
\nu\nabla^2 W = \nu D_*D\,W = \frac{1}{\rho}\!\left(\frac{\partial p}{\partial z}\right)_{\!0}.
\end{equation}

The solution of equation~\eqref{eq:8-150} is given by
\begin{equation}
\label{eq:8-152}
V = Ar + B/r,
\end{equation}
where $A$ and $B$ are determined by the speeds of rotation of the inner and the outer cylinders; they are given by equations~VII\,(145) and (146). The corresponding solution of equation~\eqref{eq:8-151} is
\begin{equation}
\label{eq:8-153}
W(r) = \frac{1}{4\rho\nu}\!\left(\frac{\partial p}{\partial z}\right)_{\!0}(r^2 + C\log r + D),
\end{equation}
where $C$ and $D$ are constants. The requirement that $W(r)$ vanishes at $r = R_1$ and $R_2$ determines $C$ and $D$ and makes the solution determinate. We find
\begin{equation}
\label{eq:8-154}
W(r) = -\frac{1}{4\rho\nu}\!\left(\frac{\partial p}{\partial z}\right)_{\!0}\biggl\{R_1^2 - r^2 + \frac{2R_1\,d + d^2}{\log(1+d/R_1)}\log\!\left(\frac{r}{R_1}\right)\biggr\},
\end{equation}
where $d\;(= R_2 - R_1)$ is the gap width.

For small gap widths ($d \ll \frac{1}{2}(R_2+R_1)$), the velocity distribution given by equation~\eqref{eq:8-154} reduces to the familiar form of the Poiseuille flow between parallel plates. Thus, letting
\begin{equation}
\label{eq:8-155}
\zeta = (r-R_1)/(R_2-R_1),
\end{equation}
we find that, to order $(d/R_1)^2$, the velocity distribution becomes
\begin{equation}
\label{eq:8-156}
W(\zeta) = 6V_m\,\zeta(1-\zeta),
\end{equation}
where
\begin{equation}
\label{eq:8-157}
V_m = -\frac{d^2}{12\rho\nu}\!\left(\frac{\partial p}{\partial z}\right)_{\!0}
\end{equation}
is the mean axial flow.

\subsection{The perturbation equations}
\label{sec:79a}

We shall now investigate the stability of the flow defined by equations~\eqref{eq:8-149}, \eqref{eq:8-152}, and \eqref{eq:8-153}. Letting
\begin{equation}
\label{eq:8-158}
u_r,\quad V + u_\theta,\quad W + u_z, \quad \text{and}\quad \varpi\;(= \delta p/\rho)
\end{equation}
describe the perturbed state, and assuming, further, that the perturbations are axisymmetric, we readily obtain from equations~VII\,(134)--(136) the linearized equations (cf.\ equations~\eqref{eq:8-69}--\eqref{eq:8-71} and equations~VII\,(149)--(153))
\begin{equation}
\label{eq:8-159}
\frac{\partial u_r}{\partial t} + W\frac{\partial u_r}{\partial z} - 2\Omega u_\theta = -\frac{\partial\varpi}{\partial r} + \nu\!\left(\nabla^2 u_r - \frac{u_r}{r^2}\right),
\end{equation}
\begin{equation}
\label{eq:8-160}
\frac{\partial u_\theta}{\partial t} + W\frac{\partial u_\theta}{\partial z} + \left(\frac{dV}{dr} + \frac{V}{r}\right)u_r = \nu\!\left(\nabla^2 u_\theta - \frac{u_\theta}{r^2}\right),
\end{equation}
\begin{equation}
\label{eq:8-161}
\frac{\partial u_z}{\partial t} + W\frac{\partial u_z}{\partial z} + u_r\frac{dW}{dr} = \nu\nabla^2 u_z,
\end{equation}
where $\nabla^2$ has now the meaning
\begin{equation}
\label{eq:8-162}
\nabla^2 = \frac{\partial^2}{\partial r^2} + \frac{1}{r}\frac{\partial}{\partial r} + \frac{\partial^2}{\partial z^2},
\end{equation}
and we also have the equation of continuity~\eqref{eq:8-72}.

Seeking solutions of equations~\eqref{eq:8-159}--\eqref{eq:8-161} which have the $(z,t)$-dependence given by \eqref{eq:8-73}, we obtain the equations
\begin{equation}
\label{eq:8-163}
\nu\!\left(DD_* - k^2 - i\frac{p}{v} - i\frac{k}{\nu}W\right)u + 2\Omega v = D\varpi,
\end{equation}
\begin{equation}
\label{eq:8-164}
\nu\!\left(DD_* - k^2 - i\frac{p}{\nu} - i\frac{k}{\nu}W\right)v = 2Au,
\end{equation}
\begin{equation}
\label{eq:8-165}
\nu\!\left(D_*D - k^2 - i\frac{p}{\nu} - i\frac{k}{\nu}W\right)w - uDW = ik\varpi,
\end{equation}
and
\begin{equation}
\label{eq:8-166}
D_*\,\mathbf{u} = -ikw,
\end{equation}
where $D$, $D_*$, $\mathbf{u}$, $v$, $w$ have the same meanings as in equations~\eqref{eq:8-74}--\eqref{eq:8-78}.

Eliminating $w$ between equations~\eqref{eq:8-165} and \eqref{eq:8-166}, we have
\begin{equation}
\label{eq:8-167}
\frac{\nu}{k}\biggl[D_*D - k^2 - i\frac{p}{\nu} - i\frac{k}{\nu}W\biggr]D_*\,\mathbf{u} + \frac{i}{k}(DW)u_r = \varpi.
\end{equation}
Inserting this expression for $\varpi$ in equation~\eqref{eq:8-163}, we obtain
\begin{multline}
\label{eq:8-168}
\frac{\nu}{k}D\!\left(D_*D - k^2 - i\frac{p}{\nu} - i\frac{k}{\nu}W\right)D_*\,\mathbf{u} + \frac{i}{k}D(uDW) \\
= \nu\!\left(DD_* - k^2 - i\frac{p}{\nu} - i\frac{k}{\nu}W\right)u + 2\Omega v.
\end{multline}
After some rearranging, equation~\eqref{eq:8-168} becomes
\begin{equation}
\label{eq:8-169}
\biggl[DD_* - k^2 - \frac{i}{k}(p+kW)\biggr]\!\biggl(DD_* - k^2\biggr)u + \frac{i}{k}\nu\biggl[DD_*\frac{W}{r}\biggr] = -\frac{2\Omega k}{\nu}\,v.
\end{equation}
This equation must be considered together with the equation
\begin{equation}
\label{eq:8-170}
\biggl[DD_* - k^2 - \frac{i}{k}(p+kW)\biggr]v = \frac{2A}{\nu}\,u.
\end{equation}
Solutions of equations~\eqref{eq:8-169} and \eqref{eq:8-170} must be sought which satisfy the usual boundary conditions
\begin{equation}
\label{eq:8-171}
\mathbf{u} = v = 0 \quad \text{and} \quad Du = 0 \quad \text{for } r = R_1 \text{ and } R_2.
\end{equation}

\subsection{The reduction to the case of a narrow gap}
\label{sec:79b}

If the gap $d = R_2 - R_1$ between the two cylinders is small compared to their mean radius $\frac{1}{2}(R_2+R_1)$, we need not, as in similar past instances, distinguish between $D$ and $D_*$, and use for $\Omega$ the linear representation
\begin{equation}
\label{eq:8-172}
\Omega = \Omega_1[1-(1-\mu)\xi].
\end{equation}
[In equation~\eqref{eq:8-172} $\xi$ is the variable defined in equation~\eqref{eq:8-155} and $\Omega_1$ and $\Omega_2\;(= \mu\Omega_1)$ are the angular velocities of the two cylinders.] In the same approximation (cf.\ equation~\eqref{eq:8-156}),
\begin{equation}
\label{eq:8-173}
Wd/\nu = 6V_m d/\nu\;\xi(1-\xi) = 6\,\mathrm{R}\,\xi(1-\xi),
\end{equation}
where R is the Reynolds number defined with respect to the mean axial flow. In the framework of these approximations which are appropriate for a narrow gap, equations~\eqref{eq:8-169} and \eqref{eq:8-170} take the forms
\begin{equation}
\label{eq:8-174}
\bigl\{(D^2-a^2) - i[\sigma + 6\,\mathrm{R}\,a\,\xi(1-\xi)]\bigr\}(D^2 - a^2)u - 12i\,\mathrm{R}\,a\,u = \frac{2\Omega_1\,d^2}{\nu}a^2[1-(1-\mu)\xi]v
\end{equation}
and
\begin{equation}
\label{eq:8-175}
\bigl\{(D^2-a^2) - i[\sigma + 6\,\mathrm{R}\,a\,\xi(1-\xi)]\bigr\}v = -\frac{2Ad^2}{\nu}\,u,
\end{equation}
where
\begin{equation}
\label{eq:8-176}
a = kd, \quad \sigma = pd^2/\nu,
\end{equation}
and $D$ stands for $d/d\xi$. By translating the origin of~$\xi$ to be midway between the two cylinders and replacing
\begin{equation}
\label{eq:8-177}
\mathbf{u} \text{ by } \tfrac{1}{2}(1+\mu)\frac{2\Omega_1\,d^2}{\nu}\,a^2 u,
\end{equation}
the equations become
\begin{equation}
\label{eq:8-178}
\bigl\{(D^2 - a^2) - i[\sigma + 6\,\mathrm{R}\,a(\tfrac{1}{4}-\xi^2)]\bigr\}(D^2 - a^2)u - 12i\,\mathrm{R}\,a\,u = \left[1 - 2\frac{1-\mu}{1+\mu}\,\xi\right]v
\end{equation}
and
\begin{equation}
\label{eq:8-179}
\bigl\{(D^2 - a^2) - i[\sigma + 6\,\mathrm{R}\,a(\tfrac{1}{4}-\xi^2)]\bigr\}v = -Ta^2 u,
\end{equation}
where
\begin{equation}
\label{eq:8-180}
T = -\tfrac{1}{2}(1+\mu)\frac{4A\Omega_1}{\nu^2}\,d^4
\end{equation}
is the Taylor number as defined in equation~VII\,(278).

Together with the boundary conditions,
\begin{equation}
\label{eq:8-181}
\mathbf{u} = Du = v = 0 \quad \text{for } \xi = \pm\tfrac{1}{2},
\end{equation}
equations~\eqref{eq:8-178} and \eqref{eq:8-179} constitute a characteristic value problem for $T$ for assigned $a$, $\mathrm{R}$, and $\sigma$. For an arbitrarily assigned $\sigma$, the characteristic values of $T$ will in general be complex. The requirement that, for a given $a$ and R, $T$ be real determines~$\sigma$. The critical Taylor number for the onset of instability (for a given R) is then given by the minimum (with respect to~$a$) of the smallest real characteristic values, $T$, so determined.

\subsection{An approximate solution of the characteristic value problem for the case $\mu > 0$}
\label{sec:79c}

In Chapter~VII, \S\,71\,(d) we saw that the critical Taylor numbers for $\mu > 0$ (and $\mathrm{R} = 0$) can be determined with fair precision by ignoring the variation of $\Omega$ across the gap and replacing $\Omega(\xi)$ by its average value. In \S\,77\,(c) we found the same thing with respect to the solution of equations~\eqref{eq:8-50} and \eqref{eq:8-51} in case $-1 \leqslant \lambda \leqslant +1$. It would, therefore, appear that for $\mu > 0$ and for Reynolds numbers R not too large, we can obtain the solutions of equations~\eqref{eq:8-178} and \eqref{eq:8-179} with comparable precision by replacing the terms in these equations, which allow for the variation of $W(\xi)$ and $\Omega(\xi)$ across the gap, by their average values. Thus, we are led to consider the simpler equations
\begin{equation}
\label{eq:8-182}
\bigl[(D^2-a^2) - i(\sigma + \mathrm{R}\,a)\bigr](D^2 - a^2)u - 12i\,\mathrm{R}\,a\,u = v
\end{equation}
and
\begin{equation}
\label{eq:8-183}
\bigl[(D^2-a^2) - i(\sigma + \mathrm{R}\,a)\bigr]v = -Ta^2 u
\end{equation}
with the same boundary conditions~\eqref{eq:8-181}.

The characteristic value problem presented by equations~\eqref{eq:8-181}--\eqref{eq:8-183} can be solved quite readily by methods with which we are now familiar; in particular, the method which was used in \S\,44\,(b) for solving equations~IV\,(141) and (142) is applicable with only minor changes. Thus, we expand $v$ in a cosine series in the form
\begin{equation}
\label{eq:8-184}
v = \sum_{m=0}^{\infty} A_m\cos[(2m+1)\pi\xi]
\end{equation}
and express $\mathbf{u}$ in the manner
\begin{equation}
\label{eq:8-185}
\mathbf{u} = \sum_{m=0}^{\infty} A_m\,\mathbf{u}_m,
\end{equation}
where $\mathbf{u}_m$ is a solution of the equation
\begin{equation}
\label{eq:8-186}
\bigl[(D^2-a^2) - i(\sigma + \mathrm{R}\,a)\bigr](D^2-a^2)u_m - 12i\,\mathrm{R}\,a\,u_m = \cos[(2m+1)\pi\xi]
\end{equation}
which satisfies the boundary conditions
\begin{equation}
\label{eq:8-187}
\mathbf{u}_m = Du_m = 0 \quad \text{for } \xi = \pm\tfrac{1}{2}.
\end{equation}
The required solution is (cf.\ equation~IV\,(175))
\begin{equation}
\label{eq:8-188}
\mathbf{u}_m = \gamma_{2m+1}\cos[(2m+1)\pi\xi] + \sum_{j=1}^{2} B_j^{(m)}\cosh q_j\,\xi,
\end{equation}
where
\begin{equation}
\label{eq:8-189}
\frac{1}{\gamma_{2m+1}} = \bigl[(2m+1)^2\pi^2 + a^2 + i(\sigma + \mathrm{R}\,a)\bigr]\bigl[(2m+1)^2\pi^2 + a^2\bigr] - 12i\,\mathrm{R}\,a,
\end{equation}
the $B_j^{(m)}$'s ($j = 1,2$) are constants of integration, and the $q_j$'s ($j = 1,2$) are the roots of the quadratic
\begin{equation}
\label{eq:8-190}
(q^2 - a^2)\bigl[(q^2-a^2) - i(\sigma + \mathrm{R}\,a)\bigr] - 12i\,\mathrm{R}\,a = 0.
\end{equation}
The constants $B_j^{(m)}$ in the solution~\eqref{eq:8-188} for $u_m$ are determined by the boundary conditions~\eqref{eq:8-187}; they are given by the same equations~IV\,(181) and (182) if we ascribe to the $\gamma$'s and the $q$'s their present meanings.

Substituting for $v$ and $\mathbf{u}$ in accordance with equations~\eqref{eq:8-184}, \eqref{eq:8-185}, and \eqref{eq:8-188} in equation~\eqref{eq:8-183}, we obtain
\begin{equation}
\label{eq:8-191}
\sum_m A_m\,C_{2m+1}\cos[(2m+1)\pi\xi] = Ta^2\sum_m A_m\bigl\{\gamma_{2m+1}\cos[(2m+1)\pi\xi] + \sum_{j=1}^{2} B_j^{(m)}\cosh q_j\,\xi\bigr\},
\end{equation}
where
\begin{equation}
\label{eq:8-192}
C_{2m+1} = (2m+1)^2\pi^2 + a^2 + i(\sigma + \mathrm{R}\,a).
\end{equation}

Equation~\eqref{eq:8-191} leads to the secular equation (cf.\ equations~IV\,(182), \eqref{eq:8-187}, and \eqref{eq:8-188})
\begin{equation}
\label{eq:8-193}
\biggl|\biggl|\frac{C_{2m+1}}{2}\gamma_{2m+1} - \gamma_{2m+1}\biggr|\biggr|\frac{8}{a_m} - (n|m) = 0,
\end{equation}
where
\begin{equation}
\label{eq:8-194}
(n|m) = (-1)^{m+n+1}\,2(2n+1)(2m+1)\pi^2\gamma_{2m+1}\Delta(q_1^2 - q_2^2)
\end{equation}
and
\begin{equation}
\label{eq:8-195}
\Delta = (q_1\tanh\tfrac{1}{2}q_1 - q_2\tanh\tfrac{1}{2}q_2)^{-1}.
\end{equation}

In the first approximation (in which we retain only the $(0,0)$-element in the secular matrix), the solution for $T$ is given by
\begin{equation}
\label{eq:8-196}
T = \frac{\pi^2 + a^2 + i(\sigma + \mathrm{R}\,a)}{a^2\gamma_1\bigl[1-4\pi^2\gamma_1\Delta(q_1^2-q_2^2)\bigr]},
\end{equation}
where
\begin{equation}
\label{eq:8-197}
\frac{1}{\gamma_1} = \bigl[\pi^2 + a^2 + i(\sigma + \mathrm{R}\,a)\bigr](\pi^2+a^2) - 12i\,\mathrm{R}\,a.
\end{equation}

\begin{table}[ht]
\centering
\caption{The critical Taylor numbers and related constants for a few Reynolds numbers}
\label{tab:XL}
\begin{tabular}{cccc}
\hline
R & $\sigma$ & $T_c$ & $-\sigma$ \\
\hline
0 & 3.12 & 1,715 & 0 \\
5 & 3.1 & 1,733 & \\
10 & 3.1 & 1,793 & 55.7 \\
20 & 3.4 & 2,399 & \\
40 & 4.2 & 3,881 & 149.9 \\
60 & 5.2 & 5,664 & 379.2 \\
80 & 6.0 & 8,319 & 475 \\
100 & 6.6 & 10,876 & 594 \\
\hline
\end{tabular}
\end{table}

The critical Taylor numbers have been determined for a few values of the Reynolds number R by making use of the solution~\eqref{eq:8-196}. The results are summarized in Table~XL; they are further illustrated in Figs.~97 and~98.

\figurePlaceholder{97}{A comparison of the observed and the theoretical dependence of the critical Taylor number $T_c$ for the onset of instability as a function of the Reynolds number R of the axial flow. The full-line curve (passing through the solid circles) represents the theoretically derived relation; and the open circles represent the experimentally derived results of Donnelly and Fultz.}

\figurePlaceholder{98}{The variation of the critical wave number $a_c$ as a function of the Reynolds number R of the mean axial flow.}

\subsection{Comparison with experimental results}
\label{sec:79d}

Several experiments reported in the literature bear on the effects of axial flow on the stability of transverse flow between rotating cylinders. Those experiments were not, however, especially designed for the study and verification of the basic hydrodynamic phenomenon; we shall, therefore, limit the present discussion to the experiments of Donnelly and Fultz which bear directly on the theoretical results of the preceding section.

The apparatus used in these experiments of Donnelly and Fultz was a modification of the one used in their earlier experiments described in \S\,74\,(b). The outer cylinder was the same precision-bore Pyrex tube of length 94~cm and of radius $R_2 = 6.2846\pm 0.006$~cm. The inner cylinder (which in the earlier experiments provided the `wide gap' with a ratio $\eta = \frac{1}{2}$) was replaced by an aluminium tube of length 90~cm and of radius $R_1 = 5.9682\pm 0.0006$~cm (over the lower effective half of the tube). The gap between the cylinders was 0.3164~cm (corresponding to $\eta = 0.9497$); the uncertainty in the gap width was about 2 per cent (implying an uncertainty from this cause of 8~per cent in the deduced Taylor numbers).

The axial flow was by gravity from above, down the annulus between the two cylinders. The water leaving the tube at the bottom made a circuit passing in succession, a gear pump, a coil in a water bath controlling the temperature, a precision thermometer, a ball-type flow meter, and finally returning to a trough at the top of the cylinder providing the axial flow.

The apparatus could be successfully operated at Reynolds numbers of the mean axial flow up to 100 and Taylor numbers in the range $10^3$ and $2.5\times 10^4$.

With zero axial flow, the period of rotation at which the onset of instability was observed was just above 7~sec corresponding to a critical Taylor number of 1700 (in good agreement with the theoretical value of 1708). Donnelly and Fultz found that the cells which appeared at marginal stability were beautifully regular and of the predicted spacing.

With a mean axial flow with $\mathrm{R}\sim 3$, Donnelly and Fultz observed the appearance of very regular cells at Taylor numbers in the range 1800--1900. Also, overstable oscillations were observed as predicted: they were manifested either in the thinning and thickening of alternate ink ribbons out of phase, or in the periodic narrowing and widening of the spacing between adjacent cell walls. The periods of oscillation for $\mathrm{R}\sim 3$ were estimated to be between 7 and 10~seconds (corresponding to $\sigma$ in the range 9 to 5). For $\mathrm{R}\sim 95$, it was definitely established that $T_c\sim 10{,}000$ and that the periods of oscillation were considerably shorter: 2.8~sec at $\mathrm{R}\sim 95$ in contrast to the 7--10~sec for $\mathrm{R}\sim 3$.

It appears then that while the experiments are still of a preliminary character, they very definitely support the theoretical expectations.

\section*{Bibliographical Notes}
\label{sec:8-bib}

\noindent\S\,76. The stability of viscous flow in a curved channel was first considered by Dean:

\begin{enumerate}
\item W.~R.~\textsc{Dean}, `Fluid motion in a curved channel', \emph{Proc.\ Roy.\ Soc.\ (London)} A, \textbf{121}, 402--20 (1928).
\end{enumerate}

\noindent Later investigations of the same problem are those of:

\begin{enumerate}\setcounter{enumi}{1}
\item W.~H.~\textsc{Reid}, `On the stability of viscous flow in a curved channel', ibid.\ \textbf{244}, 186--98 (1958).
\item G.~\textsc{Hammerlin}, `Die Stabilit\"at der Str\"omung in einem gekr\"ummten Kanal', \emph{Arch.\ Rat.\ Mech.\ and Anal.}\ \textbf{1}, 212--24 (1958).
\end{enumerate}

\noindent The treatment followed in the text is equivalent to one of two methods of solution described by Reid. Reid's other method is based on an expansion of $u$ in terms of orthogonal functions which satisfy all four of the boundary conditions on $u$. Hammerlin's treatment of the problem is based on an analytical approach which is very different from the ones adopted in the book (which are, essentially, elementary). Hammerlin's treatment derives from his earlier paper:

\begin{enumerate}\setcounter{enumi}{3}
\item G.~\textsc{Hammerlin}, `\"Uber das Eigenwertproblem der dreidimensionalen Instabilit\"at laminarer Grenzschichten an konkaven W\"anden', \emph{J.\ Rat.\ Mech.\ Analysis}, \textbf{4}, 279--321 (1955).
\end{enumerate}

\noindent Applications of similar methods to Couette flow will be found in:

\begin{enumerate}\setcounter{enumi}{4}
\item H.~\textsc{Witting}, `\"Uber den Einfluss der Stromlinienkr\"ummung auf die Stabilit\"at laminarer Str\"omungen', \emph{Arch.\ Rat.\ Mech.\ and Anal.}\ \textbf{2}, 243--83 (1958).
\end{enumerate}

\noindent The reference to the experimental work by Brewster, Grosberg, and Nissan, mentioned in \S\,(c), is given below (reference~12). Some earlier experiments on the stability of the flow in a curved channel are those of:

\begin{enumerate}\setcounter{enumi}{5}
\item M.~\textsc{Adler}, `Str\"omung in gekr\"ummten Rohren', \emph{Z.\ angew.\ Math.\ Mech.}\ \textbf{14}, 257--76 (1934).
\end{enumerate}

\noindent But Adler's experiments were not designed to determine the critical value of $\Lambda$ with precision.

\noindent There are other important aspects of flow in a curved channel, and on curved surfaces generally, which we have not considered: they relate to boundary layer phenomena which are outside the scope of this book. But reference may be made to the basic papers of:

\begin{enumerate}\setcounter{enumi}{6}
\item H.~\textsc{G\"ortler}, `\"Uber den Einfluss der Wandkr\"ummung auf die Entstehung der Turbulenz', \emph{Z.\ angew.\ Math.\ Mech.}\ \textbf{20}, 138--47 (1940).

\item ------\quad `\"Uber eine dreidimensionale Instabilit\"at laminarer Grenzschichten an konkaven W\"anden', \emph{Nachr.\ Ges.\ Wiss.\ G\"ottingen}, N.F.\ \textbf{2}, 1--26 (1940).

\item ------\quad `Instabilit\"at laminarer Grenzschichten an konkaven W\"anden gegen\"uber gewissen dreidimensionalen St\"orungen', \emph{Z.\ angew.\ Math.\ Mech.}\ \textbf{21}, 250--2 (1941).
\end{enumerate}

\noindent A general account of these investigations will be found in:

\begin{enumerate}\setcounter{enumi}{9}
\item H.~\textsc{Schlichting}, \emph{Boundary Layer Theory}, McGraw-Hill Book Company, Inc., New York, 1955; see particularly pp.~355--65.
\end{enumerate}

\medskip
\noindent\S\,77. The theoretical treatment of the problem considered in this section is due to:

\begin{enumerate}\setcounter{enumi}{10}
\item R.~C.~\textsc{Di Prima}, `The stability of viscous flow between rotating concentric cylinders with a pressure gradient acting round the cylinders', \emph{J.\ Fluid Mech.}\ \textbf{6}, 462--8 (1959).
\end{enumerate}

\noindent The reference to the experimental work described in \S\,(d) is:

\begin{enumerate}\setcounter{enumi}{11}
\item D.~B.~\textsc{Brewster}, P.~\textsc{Grosberg}, and A.~H.~\textsc{Nissan}, `The stability of viscous flow between horizontal concentric cylinders', \emph{Proc.\ Roy.\ Soc.\ (London)} A, \textbf{251}, 76--91 (1959).
\end{enumerate}

\medskip
\noindent\S\,78. The stability of inviscid flow between rotating coaxial cylinders when an axial pressure gradient is present does not seem to have attracted much attention. This section is largely an expansion of the results summarized in:

\begin{enumerate}\setcounter{enumi}{12}
\item S.~\textsc{Chandrasekhar}, `The hydrodynamic stability of inviscid flow between coaxial cylinders', \emph{Proc.\ Nat.\ Acad.\ Sci.}\ \textbf{46}, 137--41 (1960).
\end{enumerate}

\noindent \S\,78\,(a). In considering the stability of the purely axial flow, we follow the classical discussion of the analogous problem in plane-parallel flows. The following papers by Kelvin and Rayleigh are particularly relevant:

\begin{enumerate}\setcounter{enumi}{13}
\item Lord~\textsc{Kelvin}, `On a disturbing infinity in Lord Rayleigh's solution for waves in a plane vortex stratum', 186--7, \emph{Mathematical and Physical Papers}, iv, \emph{Hydrodynamics and General Dynamics}, Cambridge, England, 1910.

\item ------\quad `Rectilineal motion of viscous fluid between two parallel planes', 321--30, ibid.

\item Lord~\textsc{Rayleigh}, `On the stability, or instability, of certain fluid motions', \emph{Scientific Papers}, i, 474--7, Cambridge, England, 1899.

\item ------\quad `On the stability, or instability, of certain fluid motions.\ II', ibid.\ iii, 17--23, Cambridge, England, 1902.

\item ------\quad `On the question of the stability of the flow of fluids', ibid.\ 575--84, Cambridge, England, 1902.

\item ------\quad `On the stability, or instability, of certain fluid motions.\ III', ibid.\ iv, 203--9, Cambridge, England, 1903.

\item ------\quad `On the stability of the laminar motion of an inviscid fluid', ibid.\ vi, 197--204, Cambridge, England, 1920.
\end{enumerate}

\noindent The theorem that a necessary condition for the existence of unstable modes is that $\Psi(r)\;(= rDD_*(W)r)$ must vanish somewhere inside $(R_1,R_2)$, is the exact analogue of the one proved in references~16 and~20 for plane-parallel flows. The remarks on the character of the stable solutions, in case $\Psi(r)$ is of the same sign throughout, parallel Rayleigh's in references~17 and~20.

In reference~18 Rayleigh briefly considers the case of pure axial flow between coaxial cylinders and shows that, for non-axisymmetric perturbations having an $e^{in\theta}$-dependence, the discriminant $\Psi(r)$ is replaced by
\[
\frac{d^2W}{dr^2} + \frac{1}{r}\frac{dW}{dr} - \frac{n^2}{r^2}.
\]

The proof, that with certain restrictions on the velocity profile the condition that $\Psi(r)$ has a zero in $(R_1, R_2)$ is both necessary and sufficient for instability, is similarly patterned after Lin's simplified version of the corresponding theorem of Tollmien in the theory of plane-parallel flows. The relevant references are:

\begin{enumerate}\setcounter{enumi}{20}
\item W.~\textsc{Tollmien}, `Asymptotische Integration der St\"orungs\-differentialgleichung ebener laminarer Str\"omungen bei hohen Reynoldsschen Zahlen', \emph{Z.\ angew.\ Math.\ Mech.}\ \textbf{25/27}, 33--50 (1947).

\item ------\quad Ibid.\ 70--83 (1947).

\item C.~C.~\textsc{Lin}, `On the stability of two-dimensional parallel flows.\ I.\ General theory', \emph{Quart.\ Appl.\ Math.}\ \textbf{3}, 117--42 (1945).

\item ------\quad `On the stability of two-dimensional parallel flows.\ II.\ Stability in an inviscid fluid', ibid.\ 218--34 (1945).

\item ------\quad `On the stability of two-dimensional parallel flows.\ III.\ Stability in a viscous fluid', ibid.\ 277--301 (1946).
\end{enumerate}

\noindent For a concise account of these investigations see:

\begin{enumerate}\setcounter{enumi}{25}
\item C.~C.~\textsc{Lin}, \emph{The Theory of Hydrodynamic Stability}, Cambridge, England, 1955; see particularly \S\S\,4.3 and 8.2;
\end{enumerate}
also, reference~10 (chapter XVI, pp.~314--23). I am grateful to Professor Lin for showing me a proof by Professor~K.~O.~Friedrichs on the existence of neutral modes (under certain circumstances) which ensure the instability of the immediately neighbouring modes. The discussion in the text following equation~\eqref{eq:8-107} is modelled after Lin's treatment in reference~26 (\S\,8.2, pp.~122--3) with the simplification of Friedrichs.

\noindent The stability of viscous flow along the axis of a cylinder (the so-called Hagen--Poiseuille flow for which $\Psi(r) = 0$) has been considered by:

\begin{enumerate}\setcounter{enumi}{26}
\item J.~\textsc{Pretsch}, `\"Uber die Stabilit\"at einer Lamellarstr\"omung in einem geraden Rohr mit kreisf\"ormigem Querschnitt', \emph{Z.\ angew.\ Math.\ Mech.}\ \textbf{21}, 204--17 (1941).
\end{enumerate}

\noindent Pretsch finds that the flow (as in the inviscid case) is stable.

\noindent \S\,78\,(b). The discussion in this section is an expansion of reference~12.

\medskip
\noindent\S\,79. The first discussion of the problem treated in this section is due to:

\begin{enumerate}\setcounter{enumi}{27}
\item S.~\textsc{Goldstein}, `The stability of viscous fluid flow between rotating cylinders', \emph{Proc.\ Camb.\ Phil.\ Soc.}\ \textbf{33}, 41--61 (1937).
\end{enumerate}

\noindent However, Goldstein's insufficient treatment of the characteristic value problem has led him into gross errors in his final results. He finds, for example, that after a small initial increase the critical Taylor number $T_c$ decreases very sharply when the Reynolds number R exceeds 25; whereas the calculations summarized in \S\,(c) would indicate that $T_c$ is a monotone increasing function of R. The results summarized in Table~XL are taken from:

\begin{enumerate}\setcounter{enumi}{28}
\item S.~\textsc{Chandrasekhar}, `The hydrodynamic stability of viscid flow between coaxial cylinders', \emph{Proc.\ Nat.\ Acad.\ Sci.}\ \textbf{46}, 141--3 (1960).
\end{enumerate}

\noindent Very similar results were obtained independently and simultaneously by:

\begin{enumerate}\setcounter{enumi}{29}
\item R.~C.~\textsc{Di Prima} (in press).
\end{enumerate}

\noindent The references to the experimental work described in \S\,(d) are:

\begin{enumerate}\setcounter{enumi}{30}
\item R.~J.~\textsc{Cornish}, `Flow of water through fine clearances with relative motion of the boundaries', \emph{Proc.\ Roy.\ Soc.\ (London)} A, \textbf{140}, 227--40 (1933).

\item C.~\textsc{Gazley}, Jr., `Heat-transfer characteristics of the rotational and axial flow between concentric cylinders', \emph{Trans.\ Amer.\ Soc.\ Mech.\ Eng.}, Paper No.\ 56-A-128, 1--16 (1956).

\item J.~\textsc{Kaye} and E.~C.~\textsc{Elgar}, `Modes of adiabatic and diabatic fluid flow in an annulus with an inner rotating cylinder', ibid.\ 57-HT-14, 1--11 (1957).

\item R.~J.~\textsc{Donnelly} and \textsc{Dave Fultz}, `Experiments on the stability of spiral flow between rotating cylinders', \emph{Proc.\ Nat.\ Acad.\ Sci.}\ \textbf{46}, 1150--54 (1960).
\end{enumerate}
